% ====================================================================
\chapter{Diseño del Sistema}
% ====================================================================

\section{Arquitectura del sistema}

El diseño adopta una arquitectura híbrida. La extracción de puntos y la clasificación gestual corren en el equipo del usuario. El módulo semántico corre como servicio remoto.

\subsection{Justificación de la partición y modelo de privacidad}

Se probó una primera configuración con las tres etapas en el cliente: extracción, clasificación y traducción semántica. La extracción y la clasificación conviven con la videollamada. La demanda de memoria RAM del modelo de lenguaje obliga a retirarlo del equipo del usuario y a publicarlo como servicio en la nube. Se utiliza un filtro para evitar repeticiones ilógicas, y hacia el servicio sale solo el texto del turno en curso, ya cerrado, en forma de glosas o de oración. El video y los puntos articulares permanecen en el equipo.


\subsection{Flujo integral de datos entre componentes.}

La dinámica operativa del sistema conecta el navegador, el motor local y dos servicios remotos. La Figura~\ref{fig:diagrama_arquitectura} fija esa partición. El video y los puntos articulares quedan en el equipo. Hacia la nube sale texto: glosas en el canal directo, u oración en el canal inverso. El servicio de poses no traduce. Devuelve el archivo de la glosa.

\begin{figure}[!ht]
\centering
\resizebox{1\textwidth}{!}{%
\begin{circuitikz}
\tikzstyle{every node}=[font=\fontsize{8.0pt}{10.4pt}\selectfont]

% Caja Cliente
\draw (1.125,14.25) rectangle (5.375,7.25);
\node[anchor=north, align=center, inner xsep=0.080cm, inner ysep=0.085cm, rounded corners=0.020cm, font=\fontsize{11.9pt}{15.5pt}\selectfont] at (3.25,14.1) {Cliente};

\draw (1.875,13.25) rectangle (4.375,12.5);
\node[anchor=north, align=center, inner xsep=0.080cm, inner ysep=0.085cm, rounded corners=0.020cm, font=\fontsize{10.8pt}{14.1pt}\selectfont] at (3.125,13.1) {Google Meet};

\draw (1.875,11.375) rectangle (4.375,10.625);
\node[anchor=north, align=center, inner xsep=0.080cm, inner ysep=0.085cm, rounded corners=0.020cm, font=\fontsize{10.8pt}{14.1pt}\selectfont] at (3.125,11.225) {Extension};

\draw (1.75,8.875) rectangle (4.5,7.75);
\node[anchor=center, align=center, inner xsep=0.080cm, inner ysep=0.085cm, rounded corners=0.020cm, font=\fontsize{10.8pt}{14.1pt}\selectfont] at (3.125,8.3125) {Motor ILSA. \\ Clasificador};

% Flechas Cliente
\draw [-{Stealth[scale=1.5]}] (3.25,12.5) -- (3.25,11.375) node[pos=0.45, fill=white, fill opacity=1, text opacity=1, inner xsep=0.080cm, inner ysep=0.085cm, rounded corners=0.020cm]{fotogramas};

\draw [-{Stealth[scale=1.5]}] (2.375,10.625) -- (2.375,8.875) node[pos=0.5, align=center, fill=white, fill opacity=1, text opacity=1, inner xsep=0.080cm, inner ysep=0.085cm, rounded corners=0.020cm]{frames \\ coordenadas \\ audio};

\draw [-{Stealth[scale=1.5]}] (3.75,8.875) -- (3.75,10.625) node[pos=0.5, fill=white, fill opacity=1, text opacity=1, inner xsep=0.080cm, inner ysep=0.085cm, rounded corners=0.020cm]{subtitulos};

% Caja Nube
\draw (7.875,11.5) rectangle (11.5,7.25);
\node[anchor=north, align=center, inner xsep=0.080cm, inner ysep=0.085cm, rounded corners=0.020cm, font=\fontsize{10.8pt}{14.1pt}\selectfont] at (9.6875,11.35) {Nube};

\draw (8.375,10.5) rectangle (11,9.25);
\node[anchor=north, align=center, inner xsep=0.080cm, inner ysep=0.085cm, rounded corners=0.020cm, font=\fontsize{10.8pt}{14.1pt}\selectfont] at (9.6875,10.35) {Servicio \\ semantico};

\draw (8.375,8.875) rectangle (11,7.625);
\node[anchor=north, align=center, inner xsep=0.080cm, inner ysep=0.085cm, rounded corners=0.020cm, font=\fontsize{10.8pt}{14.1pt}\selectfont] at (9.6875,8.725) {Servicio de \\ poses};

% Flechas Cliente <-> Nube
\draw [-{Stealth[scale=1.5]}] (4.5,8.25) -- (8.375,8.25) node[pos=0.5, fill=white, fill opacity=1, text opacity=1, inner xsep=0.080cm, inner ysep=0.085cm, rounded corners=0.020cm]{nombre de glosa};

\draw [-{Stealth[scale=1.5]}] (8.375,7.875) -- (4.5,7.875) node[pos=0.5, fill=white, fill opacity=1, text opacity=1, inner xsep=0.080cm, inner ysep=0.085cm, rounded corners=0.020cm]{archivo pose};

\draw [{Stealth[scale=1.5]}-{Stealth[scale=1.5]}] (4.5,8.625) -- (8.375,9.875) node[pos=0.5, fill=white, fill opacity=1, text opacity=1, inner xsep=0.080cm, inner ysep=0.085cm, rounded corners=0.020cm]{glosas / oración};

\end{circuitikz}%
}
\caption{Diagrama de la arquitectura del sistema}
\label{fig:diagrama_arquitectura}
\end{figure}

Para modelar la interacción temporal y la secuencia de mensajes se analizan a continuación los dos canales comunicacionales de manera independiente.

\subsubsection{Dinámica temporal del Canal Señante $\rightarrow$ Oyente (Modo Sordo)}
En la Figura~\ref{fig:diagrama_secuencia_directo} se modela el ciclo de vida de un enunciado en LSA.

\begin{figure}[!ht]
\centering
\resizebox{1\textwidth}{!}{%
\begin{circuitikz}
\tikzstyle{every node}=[font=\fontsize{8.0pt}{10.4pt}\selectfont]

% Google Meet (Caja superior e inferior)
\draw (-2.25,16) rectangle (-0.625,15.125);
\node[anchor=center, align=center, inner xsep=0.080cm, inner ysep=0.085cm, rounded corners=0.020cm, font=\fontsize{8.0pt}{10.4pt}\selectfont] at (-1.4375,15.5625) {Google \\ Meet};
\draw [short] (-1.5,15.125) -- (-1.5,3.125);

\draw (-2.25,3.125) rectangle (-0.625,2.25);
\node[anchor=center, align=center, inner xsep=0.080cm, inner ysep=0.085cm, rounded corners=0.020cm, font=\fontsize{8.0pt}{10.4pt}\selectfont] at (-1.4375,2.6875) {Google \\ Meet};

% Extensión (Caja superior e inferior)
\draw (0.625,16) rectangle (2.25,15.125);
\node[anchor=center, align=center, inner xsep=0.080cm, inner ysep=0.085cm, rounded corners=0.020cm, font=\fontsize{8.0pt}{10.4pt}\selectfont] at (1.4375,15.5625) {Extensión};
\draw [short] (1.375,15.125) -- (1.375,3.125);

\draw (0.625,3.125) rectangle (2.25,2.25);
\node[anchor=center, align=center, inner xsep=0.080cm, inner ysep=0.085cm, rounded corners=0.020cm, font=\fontsize{8.0pt}{10.4pt}\selectfont] at (1.4375,2.6875) {Extensión};

% ILSA (Caja superior e inferior)
\draw (3.5,16) rectangle (5.125,15.125);
\node[anchor=center, align=center, inner xsep=0.080cm, inner ysep=0.085cm, rounded corners=0.020cm, font=\fontsize{8.0pt}{10.4pt}\selectfont] at (4.3125,15.5625) {ILSA};
\draw [short] (4.25,15.125) -- (4.25,3.125) node[pos=0.49, fill=white, fill opacity=1, text opacity=1, inner xsep=0.080cm, inner ysep=0.085cm, rounded corners=0.020cm]{cierre por temporizador};

\draw (3.5,3.125) rectangle (5.125,2.25);
\node[anchor=center, align=center, inner xsep=0.080cm, inner ysep=0.085cm, rounded corners=0.020cm, font=\fontsize{8.0pt}{10.4pt}\selectfont] at (4.3125,2.6875) {ILSA};

% Servidor semántico (Caja superior e inferior)
\draw (6.375,16) rectangle (8,15.125);
\node[anchor=center, align=center, inner xsep=0.080cm, inner ysep=0.085cm, rounded corners=0.020cm, font=\fontsize{8.0pt}{10.4pt}\selectfont] at (7.1875,15.5625) {Servidor \\ semántico};
\draw [short] (7.125,15.125) -- (7.125,3.125);

\draw (6.375,3.125) rectangle (8,2.25);
\node[anchor=center, align=center, inner xsep=0.080cm, inner ysep=0.085cm, rounded corners=0.020cm, font=\fontsize{8.0pt}{10.4pt}\selectfont] at (7.1875,2.6875) {Servidor \\ semántico};

% Loop por seña
\draw [dashed] (-1.875,15) rectangle (5.75,10.5);
\node[anchor=north east, align=right, inner xsep=0.080cm, inner ysep=0.085cm, rounded corners=0.020cm, font=\fontsize{8.0pt}{10.4pt}\selectfont] at (5.6,14.85) {[loop por seña]};

% Flujos de interacción
\draw [-{Stealth[scale=1.5]}] (-1.5,14) -- (1.375,14) node[pos=0.5, above, fill=white, fill opacity=1, text opacity=1, inner xsep=0.080cm, inner ysep=0.085cm, rounded corners=0.020cm]{video interceptado};

\draw [draw opacity=0, fill={rgb,255:red,215; green,215; blue,215}, fill opacity=1] (0.25,13.75) rectangle (3,13.25);
\node[anchor=center, align=center, inner xsep=0.080cm, inner ysep=0.085cm, rounded corners=0.020cm, font=\fontsize{8.0pt}{10.4pt}\selectfont] at (1.625,13.5) {extrae pose y manos};

\draw [-{Stealth[scale=1.5]}] (1.375,12.5) -- (4.25,12.5) node[pos=0.5, above, fill=white, fill opacity=1, text opacity=1, inner xsep=0.080cm, inner ysep=0.085cm, rounded corners=0.020cm]{POST /sign};

\draw [draw opacity=0, fill={rgb,255:red,215; green,215; blue,215}, fill opacity=1] (2.875,12.25) rectangle (5.625,11.75);
\node[anchor=center, align=center, inner xsep=0.080cm, inner ysep=0.085cm, rounded corners=0.020cm, font=\fontsize{8.0pt}{10.4pt}\selectfont] at (4.25,12) {clasificador y filtro};

\draw [-{Stealth[scale=1.5]}, dash pattern=on 6pt off 2pt] (4.25,11.375) -- (1.375,11.375) node[pos=0.5, fill=white, fill opacity=1, text opacity=1, inner xsep=0.080cm, inner ysep=0.085cm, rounded corners=0.020cm]{200 OK};

\draw [draw opacity=0, fill={rgb,255:red,215; green,215; blue,215}, fill opacity=1] (2.875,11.125) rectangle (5.625,10.625);
\node[anchor=center, align=center, inner xsep=0.080cm, inner ysep=0.085cm, rounded corners=0.020cm, font=\fontsize{8.0pt}{10.4pt}\selectfont] at (4.25,10.875) {acumula glosa};

\draw [draw opacity=0, fill={rgb,255:red,215; green,215; blue,215}, fill opacity=1] (-4.625,10.25) rectangle (4.625,9.75);
\node[anchor=center, align=center, inner xsep=0.080cm, inner ysep=0.085cm, rounded corners=0.020cm, font=\fontsize{8.0pt}{10.4pt}\selectfont] at (0,10) {pausa física de inactividad};

\draw [-{Stealth[scale=1.5]}] (4.25,8.25) -- (7.125,8.25) node[pos=0.5, above, fill=white, fill opacity=1, text opacity=1, inner xsep=0.080cm, inner ysep=0.085cm, rounded corners=0.020cm]{POST /sign};

\draw [draw opacity=0, fill={rgb,255:red,215; green,215; blue,215}, fill opacity=1] (5.75,8) rectangle (8.5,7.5);
\node[anchor=center, align=center, inner xsep=0.080cm, inner ysep=0.085cm, rounded corners=0.020cm, font=\fontsize{8.0pt}{10.4pt}\selectfont] at (7.125,7.75) {interpretación};

\draw [-{Stealth[scale=1.5]}, dash pattern=on 6pt off 2pt] (7.125,6.625) -- (4.25,6.625) node[pos=0.5, above, fill=white, fill opacity=1, text opacity=1, inner xsep=0.080cm, inner ysep=0.085cm, rounded corners=0.020cm]{200 OK \{"text": "..."\}};

\draw [-{Stealth[scale=1.5]}] (4.25,6.25) -- (1.375,6.25) node[pos=0.5, above, fill=white, fill opacity=1, text opacity=1, inner xsep=0.080cm, inner ysep=0.085cm, rounded corners=0.020cm]{oración};

\draw [-{Stealth[scale=1.5]}] (1.375,5.875) -- (-1.5,5.875) node[pos=0.5, above, fill=white, fill opacity=1, text opacity=1, inner xsep=0.080cm, inner ysep=0.085cm, rounded corners=0.020cm]{oración};

\draw [-{Stealth[scale=1.5]}] (-1.5,5.125) -- (10.125,5.125) node[pos=0.5, above, fill=white, fill opacity=1, text opacity=1, inner xsep=0.080cm, inner ysep=0.085cm, rounded corners=0.020cm]{Muestras subtítulos};

\draw [draw opacity=0, fill={rgb,255:red,215; green,215; blue,215}, fill opacity=1] (2.875,4.75) rectangle (5.625,4.25);
\node[anchor=center, align=center, inner xsep=0.080cm, inner ysep=0.085cm, rounded corners=0.020cm, font=\fontsize{8.0pt}{10.4pt}\selectfont] at (4.25,4.5) {síntesis de voz};

\draw [-{Stealth[scale=1.5]}] (4.25,3.75) -- (10.125,3.75) node[pos=0.5, above, fill=white, fill opacity=1, text opacity=1, inner xsep=0.080cm, inner ysep=0.085cm, rounded corners=0.020cm]{voz};

\draw [-{Stealth[scale=1.5]}] (-4.375,14.125) -- (-1.5,14.125) node[pos=0.5, above, fill=white, fill opacity=1, text opacity=1, inner xsep=0.080cm, inner ysep=0.085cm, rounded corners=0.020cm]{movimiento};

% Usuario señante (Izquierda)
\draw [short] (-4.375,15.125) -- (-4.375,3.625);
\draw (-4.375,16.125) circle (0.125cm);
\draw [short] (-4.375,16) -- (-4.375,15.75);
\draw [short] (-4.375,15.75) -- (-4.5,15.625);
\draw [short] (-4.375,15.75) -- (-4.25,15.625);
\draw [short] (-4.5,15.875) -- (-4.25,15.875);
\node[anchor=center, align=center, inner xsep=0.080cm, inner ysep=0.085cm, rounded corners=0.020cm, font=\fontsize{6.8pt}{8.9pt}\selectfont] at (-4.3125,15.3125) {Usuario señante};

\draw (-4.375,3.25) circle (0.125cm);
\draw [short] (-4.375,3.125) -- (-4.375,2.875);
\draw [short] (-4.375,2.875) -- (-4.5,2.75);
\draw [short] (-4.375,2.875) -- (-4.25,2.75);
\draw [short] (-4.5,3) -- (-4.25,3);
\node[anchor=center, align=center, inner xsep=0.080cm, inner ysep=0.085cm, rounded corners=0.020cm, font=\fontsize{6.8pt}{8.9pt}\selectfont] at (-4.3125,2.4375) {Usuario señante};

% Usuario oyente (Derecha)
\draw [short] (10.125,15.125) -- (10.125,3.5);
\draw (10.125,16.125) circle (0.125cm);
\draw [short] (10.125,16) -- (10.125,15.75);
\draw [short] (10.125,15.75) -- (10,15.625);
\draw [short] (10.125,15.75) -- (10.25,15.625);
\draw [short] (10,15.875) -- (10.25,15.875);
\node[anchor=center, align=center, inner xsep=0.080cm, inner ysep=0.085cm, rounded corners=0.020cm, font=\fontsize{6.8pt}{8.9pt}\selectfont] at (10.1875,15.3125) {Usuario oyente};

\draw (10.125,3.125) circle (0.125cm);
\draw [short] (10.125,3) -- (10.125,2.75);
\draw [short] (10.125,2.75) -- (10,2.625);
\draw [short] (10.125,2.75) -- (10.25,2.625);
\draw [short] (10,2.875) -- (10.25,2.875);
\node[anchor=center, align=center, inner xsep=0.080cm, inner ysep=0.085cm, rounded corners=0.020cm, font=\fontsize{6.8pt}{8.9pt}\selectfont] at (10.1875,2.3125) {Usuario oyente};

\end{circuitikz}%
}
\caption{Diagrama de secuencia del flujo de traducción de Lengua de Señas a voz y subtítulos.}
\label{fig:diagrama_secuencia_directo}
\end{figure}


\subsubsection{Dinámica temporal del Canal Oyente $\rightarrow$ Señante (Modo Oyente)}
En la Figura~\ref{fig:diagrama_secuencia_inverso} se describe la secuencia temporal que transforma el mensaje oral o escrito del oyente en representaciones visuales comprensibles en LSA para el usuario sordo.

\begin{figure}[!ht]
\centering
\resizebox{1\textwidth}{!}{%
\begin{circuitikz}
\tikzstyle{every node}=[font=\fontsize{8.0pt}{10.4pt}\selectfont]

% Google Meet (Cajas superior e inferior)
\draw (-1.125,16.75) rectangle (0.5,15.875);
\node[anchor=center, align=center, inner xsep=0.080cm, inner ysep=0.085cm, rounded corners=0.020cm, font=\fontsize{8.0pt}{10.4pt}\selectfont] at (-0.3125,16.3125) {Google \\ Meet};
\draw [short] (-0.375,15.875) -- (-0.375,2.625);

\draw (-1.125,2.625) rectangle (0.5,1.75);
\node[anchor=center, align=center, inner xsep=0.080cm, inner ysep=0.085cm, rounded corners=0.020cm, font=\fontsize{8.0pt}{10.4pt}\selectfont] at (-0.3125,2.1875) {Google \\ Meet};

% Extensión (Cajas superior e inferior)
\draw (1.75,16.75) rectangle (3.375,15.875);
\node[anchor=center, align=center, inner xsep=0.080cm, inner ysep=0.085cm, rounded corners=0.020cm, font=\fontsize{8.0pt}{10.4pt}\selectfont] at (2.5625,16.3125) {Extensión};
\draw [short] (2.5,15.875) -- (2.5,2.625);

\draw (1.75,2.625) rectangle (3.375,1.75);
\node[anchor=center, align=center, inner xsep=0.080cm, inner ysep=0.085cm, rounded corners=0.020cm, font=\fontsize{8.0pt}{10.4pt}\selectfont] at (2.5625,2.1875) {Extensión};

% ILSA (Cajas superior e inferior)
\draw (4.625,16.75) rectangle (6.25,15.875);
\node[anchor=center, align=center, inner xsep=0.080cm, inner ysep=0.085cm, rounded corners=0.020cm, font=\fontsize{8.0pt}{10.4pt}\selectfont] at (5.4375,16.3125) {ILSA};
\draw [short] (5.375,15.875) -- (5.375,2.625);

\draw (4.625,2.625) rectangle (6.25,1.75);
\node[anchor=center, align=center, inner xsep=0.080cm, inner ysep=0.085cm, rounded corners=0.020cm, font=\fontsize{8.0pt}{10.4pt}\selectfont] at (5.4375,2.1875) {ILSA};

% Servidor semántico (Cajas superior e inferior)
\draw (7.5,16.75) rectangle (9.125,15.875);
\node[anchor=center, align=center, inner xsep=0.080cm, inner ysep=0.085cm, rounded corners=0.020cm, font=\fontsize{8.0pt}{10.4pt}\selectfont] at (8.3125,16.3125) {Servidor \\ semántico};
\draw [short] (8.25,15.875) -- (8.25,2.625);

\draw (7.5,2.625) rectangle (9.125,1.75);
\node[anchor=center, align=center, inner xsep=0.080cm, inner ysep=0.085cm, rounded corners=0.020cm, font=\fontsize{8.0pt}{10.4pt}\selectfont] at (8.3125,2.1875) {Servidor \\ semántico};

% Servidor de poses (Cajas superior e inferior)
\draw (10.5,16.75) rectangle (12.125,15.875);
\node[anchor=center, align=center, inner xsep=0.080cm, inner ysep=0.085cm, rounded corners=0.020cm, font=\fontsize{8.0pt}{10.4pt}\selectfont] at (11.3125,16.3125) {Servidor \\ de poses};
\draw [short] (11.25,15.875) -- (11.25,2.625);

\draw (10.5,2.625) rectangle (12.125,1.75);
\node[anchor=center, align=center, inner xsep=0.080cm, inner ysep=0.085cm, rounded corners=0.020cm, font=\fontsize{8.0pt}{10.4pt}\selectfont] at (11.3125,2.1875) {Servidor \\ de poses};

% Usuario oyente (Izquierda)
\draw [short] (-3.25,15.875) -- (-3.25,3);
\draw (-3.25,16.875) circle (0.125cm);
\draw [short] (-3.25,16.75) -- (-3.25,16.5);
\draw [short] (-3.25,16.5) -- (-3.375,16.375);
\draw [short] (-3.25,16.5) -- (-3.125,16.375);
\draw [short] (-3.375,16.625) -- (-3.125,16.625);
\node[anchor=center, align=center, inner xsep=0.080cm, inner ysep=0.085cm, rounded corners=0.020cm, font=\fontsize{6.8pt}{8.9pt}\selectfont] at (-3.1875,16.0625) {Usuario oyente};

\draw (-3.25,2.75) circle (0.125cm);
\draw [short] (-3.25,2.625) -- (-3.25,2.375);
\draw [short] (-3.25,2.375) -- (-3.375,2.25);
\draw [short] (-3.25,2.375) -- (-3.125,2.25);
\draw [short] (-3.375,2.5) -- (-3.125,2.5);
\node[anchor=center, align=center, inner xsep=0.080cm, inner ysep=0.085cm, rounded corners=0.020cm, font=\fontsize{6.8pt}{8.9pt}\selectfont] at (-3.1875,1.9375) {Usuario oyente};

% Usuario señante (Derecha)
\draw [short] (14.25,15.875) -- (14.25,2.875);
\draw (14.25,16.875) circle (0.125cm);
\draw [short] (14.25,16.75) -- (14.25,16.5);
\draw [short] (14.25,16.5) -- (14.125,16.375);
\draw [short] (14.25,16.5) -- (14.375,16.375);
\draw [short] (14.125,16.625) -- (14.375,16.625);
\node[anchor=center, align=center, inner xsep=0.080cm, inner ysep=0.085cm, rounded corners=0.020cm, font=\fontsize{6.8pt}{8.9pt}\selectfont] at (14.3125,16.0625) {Usuario señante};

\draw (14.25,2.625) circle (0.125cm);
\draw [short] (14.25,2.5) -- (14.25,2.25);
\draw [short] (14.25,2.25) -- (14.125,2.125);
\draw [short] (14.25,2.25) -- (14.375,2.125);
\draw [short] (14.125,2.375) -- (14.375,2.375);
\node[anchor=center, align=center, inner xsep=0.080cm, inner ysep=0.085cm, rounded corners=0.020cm, font=\fontsize{6.8pt}{8.9pt}\selectfont] at (14.3125,1.8125) {Usuario señante};

% Marcos de agrupación condicionales y de bucles
\draw [dash pattern=on 6pt off 2pt] (2.125,12.125) rectangle (5.75,14.5);
\node[anchor=north east, align=right, inner xsep=0.080cm, inner ysep=0.085cm, rounded corners=0.020cm, font=\fontsize{8.0pt}{10.4pt}\selectfont] at (5.6,14.35) {[Si viene de micrófono]};

\draw [dash pattern=on 6pt off 2pt] (4.875,4.375) rectangle (11.875,9.25);
\node[anchor=north east, align=right, inner xsep=0.080cm, inner ysep=0.085cm, rounded corners=0.020cm, font=\fontsize{8.0pt}{10.4pt}\selectfont] at (11.725,9.1) {[Por cada glosa]};

\draw [dash pattern=on 6pt off 2pt] (5.125,5.625) rectangle (11.5,8);
\node[anchor=north east, align=right, inner xsep=0.080cm, inner ysep=0.085cm, rounded corners=0.020cm, font=\fontsize{8.0pt}{10.4pt}\selectfont] at (11.35,7.85) {[No existe archivo]};

% Cuadro de proceso de interpretación
\draw [draw opacity=0, fill={rgb,255:red,215; green,215; blue,215}, fill opacity=1] (6.75,10.375) rectangle (9.5,10.875);
\node[anchor=center, align=center, inner xsep=0.080cm, inner ysep=0.085cm, rounded corners=0.020cm, font=\fontsize{8.0pt}{10.4pt}\selectfont] at (8.125,10.625) {interpretación};

% Flujos de mensajes (Flechas)
\draw [-{Stealth[scale=1.5]}] (-3.25,14.875) -- (-0.375,14.875) node[pos=0.5, above, fill=white, fill opacity=1, text opacity=1, inner xsep=0.080cm, inner ysep=0.085cm, rounded corners=0.020cm]{habla o escribe};

\draw [-{Stealth[scale=1.5]}] (-0.375,14.75) -- (2.5,14.75) node[pos=0.5, above, fill=white, fill opacity=1, text opacity=1, inner xsep=0.080cm, inner ysep=0.085cm, rounded corners=0.020cm]{audio/texto chat};

\draw [-{Stealth[scale=1.5]}] (2.5,13.25) -- (5.375,13.25) node[pos=0.5, above, fill=white, fill opacity=1, text opacity=1, inner xsep=0.080cm, inner ysep=0.085cm, rounded corners=0.020cm]{POST /stt};

\draw [-{Stealth[scale=1.5]}, dash pattern=on 6pt off 2pt] (5.375,12.5) -- (2.5,12.5) node[pos=0.5, above, fill=white, fill opacity=1, text opacity=1, inner xsep=0.080cm, inner ysep=0.085cm, rounded corners=0.020cm]{oración};

\draw [-{Stealth[scale=1.5]}] (2.5,11.375) -- (5.375,11.375) node[pos=0.5, above, fill=white, fill opacity=1, text opacity=1, inner xsep=0.080cm, inner ysep=0.085cm, rounded corners=0.020cm]{POST /hearing};

\draw [-{Stealth[scale=1.5]}] (5.375,11.125) -- (8.25,11.125) node[pos=0.5, above, fill=white, fill opacity=1, text opacity=1, inner xsep=0.080cm, inner ysep=0.085cm, rounded corners=0.020cm]{POST /hearing};

\draw [-{Stealth[scale=1.5]}, dash pattern=on 6pt off 2pt] (8.25,9.5) -- (5.375,9.5) node[pos=0.5, above, fill=white, fill opacity=1, text opacity=1, inner xsep=0.080cm, inner ysep=0.085cm, rounded corners=0.020cm]{200 OK, glosas};

\draw [-{Stealth[scale=1.5]}] (5.375,8.375) -- (11.25,8.375) node[pos=0.5, above, fill=white, fill opacity=1, text opacity=1, inner xsep=0.080cm, inner ysep=0.085cm, rounded corners=0.020cm]{GET /tposes/\{archivo\}};

\draw [-{Stealth[scale=1.5]}, dash pattern=on 6pt off 2pt] (11.25,6.75) -- (5.375,6.75) node[pos=0.5, above, fill=white, fill opacity=1, text opacity=1, inner xsep=0.080cm, inner ysep=0.085cm, rounded corners=0.020cm]{404 NOT FOUND};

\draw [-{Stealth[scale=1.5]}] (5.375,6) -- (11.25,6) node[pos=0.5, above, fill=white, fill opacity=1, text opacity=1, inner xsep=0.080cm, inner ysep=0.085cm, rounded corners=0.020cm]{GET /fingerspell?letters=...};

\draw [-{Stealth[scale=1.5]}, dash pattern=on 6pt off 2pt] (11.25,4.75) -- (5.375,4.75) node[pos=0.5, above, fill=white, fill opacity=1, text opacity=1, inner xsep=0.080cm, inner ysep=0.085cm, rounded corners=0.020cm]{200, archivo/s de seña/s};

\draw [-{Stealth[scale=1.5]}, dash pattern=on 6pt off 2pt] (5.375,4) -- (2.5,4) node[pos=0.5, above, fill=white, fill opacity=1, text opacity=1, inner xsep=0.080cm, inner ysep=0.085cm, rounded corners=0.020cm]{Glosas y animación};

\draw [-{Stealth[scale=1.5]}] (2.5,3.75) -- (-0.375,3.75) node[pos=0.5, above, fill=white, fill opacity=1, text opacity=1, inner xsep=0.080cm, inner ysep=0.085cm, rounded corners=0.020cm]{Subtítulo/esqueleto};

\draw [-{Stealth[scale=1.5]}] (-0.375,3.375) -- (14.25,3.375) node[pos=0.5, above, fill=white, fill opacity=1, text opacity=1, inner xsep=0.080cm, inner ysep=0.085cm, rounded corners=0.020cm]{Avatar señando / subtítulos};

\end{circuitikz}%
}
\caption{Diagrama de secuencia del flujo de interpretación de voz/texto a señas para el usuario señante.}
\label{fig:diagrama_secuencia_inverso}
\end{figure}

\section{Diseño de contratos de API REST}

La comunicación entre el cliente y los servicios se rige por contratos tipados en formato JSON.

\subsection{API local del motor ILSA}

El motor ILSA expone una interfaz HTTP en el bucle local, en \texttt{127.0.0.1:8765}. La extensión de Meet es el cliente y el proceso del motor es el servidor. El contrato cubre el modo señante, el modo oyente y dos recursos de archivo. La mayoría de las respuestas usa JSON. \texttt{GET /mediapipe/\{name\}} y \texttt{GET /download/exe} devuelven el archivo en binario.

Los métodos son \texttt{GET} y \texttt{POST}. Un cuerpo inválido, un modo incorrecto o un nombre de archivo no permitido producen \texttt{400 Bad Request}. Si el proceso todavía no tiene sesión, las rutas que consultan esa sesión responden \texttt{503 Service Unavailable}. La transcripción responde \texttt{502} cuando el reconocedor de voz falla. Un archivo ausente responde \texttt{404}. Las rutas de salud, de métricas, de transcripción, de MediaPipe y del ejecutable no exigen una sesión ya abierta.

\begin{figure}[!ht]
\centering
\resizebox{1\textwidth}{!}{%
\begin{circuitikz}
\tikzstyle{every node}=[font=\fontsize{8.0pt}{10.4pt}\selectfont]

% Contenedor principal: ILSA (motor local)
\draw [line width=0.6pt] (0.625,15) rectangle (12.625,4.875);
\node[anchor=north west, align=left, inner xsep=0.080cm, inner ysep=0.085cm, rounded corners=0.020cm, font=\fontsize{10.8pt}{14.1pt}\selectfont] at (0.775,14.85) {ILSA (motor local)};

% Caja Extension
\draw (0.5,16.875) rectangle (12.625,16.125);
\node[anchor=north, align=center, inner xsep=0.080cm, inner ysep=0.085cm, rounded corners=0.020cm, font=\fontsize{10.8pt}{14.1pt}\selectfont] at (6.5625,16.725) {Extension};

% Subcaja Clasificador
\draw [line width=0.7pt] (1,11) rectangle (4.625,5.25);
\node[anchor=north west, align=left, inner xsep=0.080cm, inner ysep=0.085cm, rounded corners=0.020cm, font=\fontsize{10.8pt}{14.1pt}\selectfont] at (1.15,10.85) {clasificador};

\draw (1.25,10) rectangle (4.375,8.75);
\node[anchor=north, align=center, inner xsep=0.080cm, inner ysep=0.085cm, rounded corners=0.020cm, font=\fontsize{10.8pt}{14.1pt}\selectfont] at (2.8125,9.85) {extraccion de \\ puntos corporales};

\draw (1.25,7.875) rectangle (4.375,7.125);
\node[anchor=north, align=center, inner xsep=0.080cm, inner ysep=0.085cm, rounded corners=0.020cm, font=\fontsize{10.8pt}{14.1pt}\selectfont] at (2.8125,7.725) {clasificador};

\draw (1.25,6.25) rectangle (4.375,5.5);
\node[anchor=north, align=center, inner xsep=0.080cm, inner ysep=0.085cm, rounded corners=0.020cm, font=\fontsize{10.8pt}{14.1pt}\selectfont] at (2.8125,6.1) {filtros y buffer};

% Cajas de síntesis y sesión
\draw (9.875,11) rectangle (12.375,9.875);
\node[anchor=center, align=center, inner xsep=0.080cm, inner ysep=0.085cm, rounded corners=0.020cm, font=\fontsize{10.8pt}{14.1pt}\selectfont] at (11.125,10.4375) {sintesis de \\ voz};

\draw (5.125,11) rectangle (9.375,9.875);
\node[anchor=center, align=center, inner xsep=0.080cm, inner ysep=0.085cm, rounded corners=0.020cm, font=\fontsize{10.8pt}{14.1pt}\selectfont] at (7.25,10.4375) {sesion oyente y avatar};

% Conexiones internas
\draw [-{Stealth[scale=1.5]}] (6,16.125) -- (6,11) node[pos=0.63, fill=white, fill opacity=1, text opacity=1, inner xsep=0.080cm, inner ysep=0.085cm, rounded corners=0.020cm]{POST /hearing};
\draw [short] (6.5,11.875) -- (6.5,11.875);
\draw [short] (4.25,12.25) -- (4.25,12.25);
\draw [short] (3.5,11.75) -- (3.5,11.75);

\draw [-{Stealth[scale=1.5]}] (2.75,8.75) -- (2.75,7.875) node[pos=0.41, fill=white, fill opacity=1, text opacity=1, inner xsep=0.080cm, inner ysep=0.085cm, rounded corners=0.020cm]{puntos corporales};
\draw [-{Stealth[scale=1.5]}] (2.75,7.125) -- (2.75,6.25) node[pos=0.38, fill=white, fill opacity=1, text opacity=1, inner xsep=0.080cm, inner ysep=0.085cm, rounded corners=0.020cm]{glosa y confianza};

% Contenedor Nube
\draw (0.625,3.75) rectangle (11.5,3);
\node[anchor=north, align=center, inner xsep=0.080cm, inner ysep=0.085cm, rounded corners=0.020cm, font=\fontsize{10.8pt}{14.1pt}\selectfont] at (6.0625,3.6) {Nube};

% Flujos con la Nube y la Extensión
\draw [-{Stealth[scale=1.5]}] (5.5,9.875) -- (5.5,3.75) node[pos=0.92, fill=white, fill opacity=1, text opacity=1, inner xsep=0.080cm, inner ysep=0.085cm, rounded corners=0.020cm]{GET /tposes/\{archivo\}};
\draw [-{Stealth[scale=1.5]}, dash pattern=on 6pt off 2pt] (6.75,3.75) -- (6.75,9.875) node[pos=0.6, fill=white, fill opacity=1, text opacity=1, inner xsep=0.080cm, inner ysep=0.085cm, rounded corners=0.020cm]{archivo pose};

\draw [-{Stealth[scale=1.5]}] (5.125,16.125) -- (5.125,13.625) -- (3.25,13.625) -- (3.25,10);
\node[fill=white, fill opacity=1, text opacity=1, inner xsep=0.080cm, inner ysep=0.085cm, rounded corners=0.020cm, font=\fontsize{8.0pt}{10.4pt}\selectfont] at (4.125,13.625) {POST /sign};

\draw [-{Stealth[scale=1.5]}] (11,3.75) -- (11,9.875) node[pos=0.5, fill=white, fill opacity=1, text opacity=1, inner xsep=0.080cm, inner ysep=0.085cm, rounded corners=0.020cm]{oración};
\draw [-{Stealth[scale=1.5]}, dash pattern=on 6pt off 2pt] (6.75,11) -- (6.75,16.125) node[pos=0.61, fill=white, fill opacity=1, text opacity=1, inner xsep=0.080cm, inner ysep=0.085cm, rounded corners=0.020cm]{animación};

\draw [-{Stealth[scale=1.5]}] (1.5,5.5) -- (1.5,3.75) node[pos=0.64, fill=white, fill opacity=1, text opacity=1, inner xsep=0.080cm, inner ysep=0.085cm, rounded corners=0.020cm]{POST /sign};
\draw [-{Stealth[scale=1.5]}, dash pattern=on 6pt off 2pt] (3.125,3.75) -- (3.125,5.5) node[pos=0.37, fill=white, fill opacity=1, text opacity=1, inner xsep=0.080cm, inner ysep=0.085cm, rounded corners=0.020cm]{oración};

\draw [dash pattern=on 6pt off 2pt, -{Stealth[scale=1.5]}] (1.25,5.75) -- (-1.375,5.75) -- (-1.375,16.5) -- (0.5,16.5);
\node[fill=white, fill opacity=1, text opacity=1, inner xsep=0.080cm, inner ysep=0.085cm, rounded corners=0.020cm, font=\fontsize{8.0pt}{10.4pt}\selectfont] at (-1.375,8.75) {oración};

\draw [-{Stealth[scale=1.5]}] (8,9.875) -- (8,3.75) node[pos=0.59, fill=white, fill opacity=1, text opacity=1, inner xsep=0.080cm, inner ysep=0.085cm, rounded corners=0.020cm]{POST /hearing};
\draw [-{Stealth[scale=1.5]}, dash pattern=on 6pt off 2pt] (9,3.75) -- (9,9.875) node[pos=0.53, fill=white, fill opacity=1, text opacity=1, inner xsep=0.080cm, inner ysep=0.085cm, rounded corners=0.020cm]{glosas};

\draw [dash pattern=on 6pt off 2pt, -{Stealth[scale=1.5]}] (11,11) -- (11,14.25) -- (11,16.125);
\node[fill=white, fill opacity=1, text opacity=1, inner xsep=0.080cm, inner ysep=0.085cm, rounded corners=0.020cm, font=\fontsize{8.0pt}{10.4pt}\selectfont] at (11,13.625) {audio};

\draw [-{Stealth[scale=1.5]}] (0.375,16.25) -- (-0.125,16.25) -- (-0.125,6) -- (1.25,6);
\node[fill=white, fill opacity=1, text opacity=1, inner xsep=0.080cm, inner ysep=0.085cm, rounded corners=0.020cm, align=center, font=\fontsize{8.0pt}{10.4pt}\selectfont] at (-0.125,12.625) {POST /session \\ POST /activity \\POST /utterance/end \\GET /health \\GET /state};

\draw [-{Stealth[scale=1.5]}] (8,16.125) -- (8,13.875) -- (8,12.375) -- (8,11);
\node[fill=white, fill opacity=1, text opacity=1, inner xsep=0.080cm, inner ysep=0.085cm, rounded corners=0.020cm, font=\fontsize{8.0pt}{10.4pt}\selectfont] at (8,12.375) {POST /stt};

\draw [dash pattern=on 6pt off 2pt, -{Stealth[scale=1.5]}] (8.875,11) -- (8.875,14.5) -- (8.875,16.125);
\node[fill=white, fill opacity=1, text opacity=1, inner xsep=0.080cm, inner ysep=0.085cm, rounded corners=0.020cm, font=\fontsize{8.0pt}{10.4pt}\selectfont] at (8.875,11.75) {oración};

\end{circuitikz}%
}
\caption{Arquitectura detallada del motor local ILSA y su interacción con los servicios en la nube.}
\label{fig:ilsa_local_arquitectura}
\end{figure}

\subsubsection{Cuadro resumen de endpoints}

El Cuadro~\ref{tab:api-ilsa-local} reúne todas las rutas del motor. El estado de las glosas reside en el búfer del proceso, no en la extensión.

\begin{table}[htbp]
  \centering
  \caption{Endpoints de la API local del motor ILSA
    (\texttt{http://127.0.0.1:8765}).}
  \label{tab:api-ilsa-local}
  \scriptsize
  \begin{tabular}{|p{0.10\textwidth}|p{0.24\textwidth}|p{0.48\textwidth}|p{0.12\textwidth}|}
    \hline
    \textbf{Método} & \textbf{Ruta} & \textbf{Contrato} & \textbf{Códigos} \\ \hline
    GET & \texttt{/health} &
    Disponibilidad del proceso, del modo, del clasificador y del traductor remoto. No altera estado. &
    200 \\ \hline
    GET & \texttt{/config} &
    Umbrales de captura del modo en curso. &
    200, 503 \\ \hline
    GET & \texttt{/state} &
    Búfer de glosas, último texto y estado del traductor. &
    200, 503 \\ \hline
    GET & \texttt{/metrics} &
    Conteos y tiempos de la sesión en curso. &
    200 \\ \hline
    POST & \texttt{/session} &
    Nueva sesión. El cuerpo fija la lateralidad con \texttt{left\_handed}. La respuesta es la misma foto que \texttt{GET /state}. &
    200, 503 \\ \hline
    POST & \texttt{/sign} &
    Solo modo señante. Recibe hasta 60 fotogramas, cada uno con pose y manos. El motor los deja en una matriz de $16 \times 225$ y devuelve si la glosa entra al búfer. &
    200, 400, 503 \\ \hline
    POST & \texttt{/stt} &
    Solo modo oyente. Recibe audio WAV en Base64 y devuelve la oración en español. &
    200, 400, 502 \\ \hline
    POST & \texttt{/hearing} &
    Solo modo oyente. Recibe la oración y una marca de cierre. Consulta el canal inverso y, si el turno cerró, pide los archivos de pose. &
    200, 400, 503 \\ \hline
    POST & \texttt{/activity} &
    Sin cuerpo. Reinicia la pausa de unos 4\,s y aplaza el cierre del enunciado. &
    200, 503 \\ \hline
    POST & \texttt{/utterance/end} &
    Cierra el enunciado del modo señante, envía las glosas al traductor y vacía el búfer. &
    200, 503 \\ \hline
    GET & \texttt{/mediapipe/\{name\}} &
    Archivo local de MediaPipe para la extensión. Admite \texttt{.wasm}, \texttt{.js}, \texttt{.data}, \texttt{.tflite} y \texttt{.binarypb}. &
    200, 400, 404 \\ \hline
    GET & \texttt{/download/exe} &
    Ejecutable del motor. Si no está en el equipo, entrega el archivo de arranque. &
    200, 404 \\ \hline
  \end{tabular}
\end{table}


\texttt{GET /health} indica si el proceso está en pie, en qué modo opera y si el clasificador y el traductor responden. \texttt{GET /state} y \texttt{POST /session} comparten la foto del motor: lateralidad, glosas pendientes, último español y estado del traductor. \texttt{POST /activity} solo mueve el reloj de pausa. \texttt{POST /utterance/end} cierra el enunciado aunque ese reloj no haya vencido.

En modo oyente, \texttt{POST /sign} responde 400. \texttt{POST /stt} transcribe el micrófono de la videollamada. \texttt{POST /hearing} traduce la oración y adjunta el clip. En modo señante, \texttt{POST /stt} y \texttt{POST /hearing} responden 400. De esas dos, solo \texttt{POST /hearing} consulta la sesión y puede responder 503.

\subsubsection{Contrato POST \texttt{/session}}

La apertura de sesión declara la lateralidad del señante. El valor \texttt{false} corresponde a ejecución diestra; \texttt{true} invierte el eje horizontal de los puntos de pose y manos antes de la clasificación. Una petición válida recibe \texttt{200 OK} y la foto del motor, la misma de \texttt{GET /state}, ya con la lateralidad aplicada.

Petición:
\begin{lstlisting}[language=json]
{
  "left_handed": false
}
\end{lstlisting}

Respuesta \texttt{200 OK} (campos principales):

\begin{lstlisting}[language=json]
{
  "left_handed": false,
  "classifier_ready": true,
  "glosses": [],
  "spanish": "",
  "semantic_ready": true
}
\end{lstlisting}

\subsubsection{Contrato POST \texttt{/sign}}

La ruta recibe una seña ya recortada, y solo en modo señante. El campo \texttt{frames} es una lista de objetos. Cada objeto trae la pose y las dos manos. El tope es 60 fotogramas. El motor suaviza esa lista y la deja en una matriz de $16 \times 225$ ($33\times 3$ de pose y $21\times 3$ por mano). Un cuerpo vacío responde \texttt{400}. En modo oyente la misma ruta responde \texttt{400} y no clasifica.

Petición (secuencia abreviada):

\begin{lstlisting}[language=json]
{
  "frames": [
    {"pose": [], "left_hand": [], "right_hand": []}
  ],
  "client": {}
}
\end{lstlisting}

Una respuesta \texttt{200} puede aceptar o rechazar la seña. El rechazo por pausa entre inferencias, por seña demasiado corta o por una matriz inválida sigue en \texttt{200}, con \texttt{accepted} en falso y un motivo. No usa el código 400. Ese código queda para el cuerpo vacío y para el modo oyente.

\begin{lstlisting}[language=json]
{
  "accepted": true,
  "added": true,
  "gloss": "HOLA",
  "top3": [
    {"gloss": "HOLA", "confidence": 0.91}
  ],
  "glosses": ["HOLA"]
}
\end{lstlisting}

\texttt{added} indica si la glosa entró al búfer. Una repetición no admitida o una confianza baja dejan \texttt{accepted} en verdadero y \texttt{added} en falso. \texttt{top3} ordena hasta tres clases. El cierre del enunciado no viaja en esta respuesta: lo decide la pausa o \texttt{POST /utterance/end}.

\subsection{API del servicio semántico remoto}
\label{sec:api-semantico-remoto}

La arquitectura local se comunica con el microservicio semántico en la nube mediante una API REST sobre HTTPS. Cada petición es independiente: el contrato no incluye memoria de turnos previos, es \textit{serverless}. El canal directo corresponde a la traducción de glosas LSA a una oración en español rioplatense. El canal inverso corresponde a la traducción de una oración en español a una secuencia de glosas con orden \textit{Topic-Comment}, apta para la animación del avatar.

El Cuadro~\ref{tab:api-semantico-endpoints} resume los endpoints del diseño.

\begin{table}[htbp]
  \centering
  \caption{Endpoints del servicio semántico remoto.}
  \label{tab:api-semantico-endpoints}
  \begin{tabular}{@{}llp{8.2cm}@{}}
    \toprule
    Método & Ruta & Función \\
    \midrule
    GET & \texttt{/health} &
    Consulta de disponibilidad: indica si los modelos de lenguaje (SLM) residen en memoria RAM y si el servicio admite inferencia. \\
    POST & \texttt{/sordo}, \texttt{/translate} &
    Canal directo (LSA $\rightarrow$ español rioplatense). Las dos rutas reciben la misma secuencia de glosas y devuelven la misma oración. \\
    POST & \texttt{/hearing} &
    Canal inverso (español $\rightarrow$ LSA Topic-Comment). Recibe el texto del interlocutor oyente y devuelve la secuencia de glosas. Esas glosas alimentan después el servicio de archivos del avatar. \\
    \bottomrule
  \end{tabular}
\end{table}

\subsubsection{Contrato de \texttt{POST /sordo} y \texttt{POST /translate} (canal directo)}

\texttt{POST /sordo} y \texttt{POST /translate} comparten este contrato. La petición transporta las glosas reconocidas en el edge como una única cadena, con tokens en mayúsculas separados por espacios. La respuesta, en caso de éxito (código 200), incluye la oración en español, el identificador del SLM y el tiempo de inferencia en milisegundos (\texttt{ms}).

Petición:

\begin{lstlisting}[language=json]
{
  "glosses": "APELLIDO MIO PIÑON"
}
\end{lstlisting}

Respuesta (200 OK):

\begin{lstlisting}[language=json]
{
  "spanish": "Mi apellido es Piñon.",
  "glosses": "APELLIDO MIO PIÑON",
  "model_used": "qwen2.5-1.5b.gguf",
  "ms": 412.3,
  "queued_ms": 0.0
}
\end{lstlisting}

\subsubsection{Contrato de \texttt{POST /hearing} (canal inverso)}

La petición transporta la oración en español rioplatense, ya transcripta en el edge. La respuesta, en caso de éxito (código 200), incluye la secuencia de glosas en una cadena de tokens en mayúsculas, el SLM empleado y el tiempo de inferencia.

Petición:

\begin{lstlisting}[language=json]
{
  "spanish": "Hola, ¿cómo estás?"
}
\end{lstlisting}

Respuesta (200 OK):

\begin{lstlisting}[language=json]
{
  "glosses": "HOLA VOS COMO",
  "spanish": "Hola, ¿cómo estás?",
  "model_used": "qwen2.5-1.5b.gguf",
  "ms": 380.1,
  "queued_ms": 0.0
}
\end{lstlisting}

El campo \texttt{queued\_ms} registra la espera en cola si el SLM atiende otra inferencia; no forma parte del texto traducido. Un código 503 indica que el modelo no está listo en RAM. El endpoint \texttt{GET /health} permite al \textit{edge} diferir el encendido del clasificador local.

\subsection{API del servicio de archivos del avatar}
\label{sec:api-avatar}

El canal inverso no anima por sí solo. Un segundo servicio remoto entrega el archivo de movimiento de cada glosa. El nombre del archivo, sin la extensión, es la glosa ya normalizada. El puerto, la imagen y la plataforma de ese servicio quedan en la Sección~\ref{sec:impl_tpose}.

El Cuadro~\ref{tab:api-avatar-endpoints} resume las rutas. La de salud responde JSON. La del archivo responde el binario, no un documento JSON.

\begin{table}[htbp]
  \centering
  \caption{Endpoints del servicio de archivos del avatar.}
  \label{tab:api-avatar-endpoints}
  \begin{tabular}{@{}llp{7.6cm}@{}}
    \toprule
    Método & Ruta & Función \\
    \midrule
    GET & \texttt{/health} &
    Disponibilidad del proceso. Responde 200 con \texttt{\{"status": "ok"\}}. \\
    GET & \texttt{/tposes/\{glosa\}.pose} &
    Archivo de esa glosa. El nombre debe coincidir de forma exacta, extensión incluida. 200 y cuerpo binario si existe. 404 si no existe. \\
    GET & \texttt{/fingerspell} &
    Solo letras del alfabeto manual, varias en el parámetro \texttt{letters}, separadas por coma. Por ejemplo, \texttt{GET /fingerspell?letters=H,O,L,A}. 200 y el archivo de cada letra. 400 si el pedido viene vacío o trae algo que no es letra. 404 si falta el archivo de una de esas letras. \\
    \bottomrule
  \end{tabular}
\end{table}

ILSA arma el nombre antes de la solicitud. El resultado es una sola cadena de letras mayúsculas y números, sin tildes, sin espacios y sin otros signos. Las reglas se aplican en este orden:
\begin{enumerate}
    \item Toda la glosa pasa a mayúsculas.
    \item Queda la forma masculina.
    \item Las vocales con tilde pasan a la vocal simple. \texttt{CAFÉ} queda \texttt{CAFE}.
    \item Se borran guiones bajos.
    \item La letra \texttt{Ñ}, cuando esa es toda la glosa, se pide como \texttt{ENIE}. Una \texttt{Ñ} dentro de la palabra pasa a \texttt{N}. \texttt{ESPAÑA} queda \texttt{ESPANA}.
    \item La glosa de parentesco \texttt{PAPA} se pide como \texttt{PADRE}. El archivo \texttt{PAPA.pose} corresponde al alimento. Para el parentesco, ILSA usa \texttt{PADRE}.
\end{enumerate}

La dirección del servicio y la caché local se describen en la Sección~\ref{sec:impl_tpose}. La búsqueda de una glosa compuesta queda en la Sección~\ref{sec:repr_inversa}.

\subsubsection{Ausencia de un archivo}
\label{sec:ausencia_archivo}

Cuando el nombre pedido no tiene archivo, el servicio puede cerrar ese hueco de cuatro maneras. Las cuatro se consideraron. El prototipo adopta la cuarta.

\paragraph{Letra por letra.}
Un archivo ausente responde 404. El cliente pide después cada letra del deletreo en una solicitud propia.
A favor, cada letra queda como un archivo independiente, y una letra que falta no anula a las demás.
En contra, un nombre largo exige un viaje de red por letra, después del 404 inicial.

\paragraph{Lote de letras o de glosas.}
El archivo ausente sigue en 404. Otra solicitud admite varias letras juntas, y también todas las glosas del enunciado.
A favor, el deletreo, o el enunciado entero, viaja en un solo pedido.
En contra, la misma solicitud mezcla la seña léxica con la letra. Si falta un nombre en el medio, la respuesta deja de ser un archivo o un 404 y tiene que explicar un éxito parcial.

\paragraph{Deletreo en la misma ruta.}
La ruta que pide la glosa, si no hay archivo, responde con las letras que forman ese nombre.
A favor, el cliente no hace un segundo pedido: el deletreo vuelve en la misma llamada.
En contra, una misma ruta a veces entrega una seña y a veces una lista de letras. El 404, que avisa que el diccionario no tiene esa glosa, desaparece. Una caché bajo ese nombre guardaría el deletreo en el lugar de una seña que todavía no existe.

\paragraph{Ruta exclusiva de deletreo.}
El archivo ausente responde 404. El deletreo tiene su propia ruta, \texttt{GET /fingerspell}, y esa ruta solo acepta letras. Varias letras entran en una sola solicitud, como en \texttt{GET /fingerspell?letters=H,O,L,A}.
A favor, la ruta del archivo conserva un nombre y un archivo, o un 404. La ruta del deletreo no acepta una glosa léxica, y un nombre no cuesta un viaje por letra.
En contra, después del 404 el cliente hace una segunda solicitud, la del deletreo.

La cuarta opción es la retenida. El 404 conserva un significado único: esa glosa no tiene seña. El cliente, y no el servicio, decide si ese nombre se deletrea. La primera opción conserva el mismo 404, pero multiplica los viajes en cada nombre propio. La segunda ahorra viajes, y a la vez abre la puerta a pedir todo el enunciado de una vez: un hueco en el medio obliga entonces a un contrato de éxito parcial. La tercera ahorra el segundo viaje, y esconde la ausencia dentro de una respuesta exitosa: quien llama ya no distingue la falta de seña de la seña misma. La ruta aparte deja el archivo como lectura exacta, apta para la caché y para las réplicas de solo lectura, y concentra el deletreo en un pedido que solo entiende letras. \texttt{CH} y \texttt{Ñ} cuentan como una letra, igual que en la Sección~\ref{sec:contingencia_lexica}.

\section{Diseño de la representación gestual inversa}
\label{sec:repr_inversa}

El canal oyente recibe del traductor inverso una lista de glosas en mayúsculas. Esa lista, por sí sola, es texto: la persona sorda todavía no ve la seña. El diseño convierte cada glosa en el movimiento correspondiente, tomado de un archivo que ya contiene la seña medida sobre un video de referencia. El archivo no guarda una malla ni las rotaciones de un muñeco. Guarda, fotograma a fotograma, la posición de puntos articulares (\textit{landmarks}): hombros, codos, muñecas, las dos manos y un contorno de cara. El servicio de la Sección~\ref{sec:api-avatar} entrega ese archivo. Lo que sigue fija cómo se prepara antes de mostrarlo, cómo se une una seña con la siguiente y qué ocurre cuando la glosa no tiene archivo. La tasa común, el peso del suavizado, el tamaño del lienzo, el umbral de confianza, la longitud del puente y la duración de la entrada y de la salida quedan en la Sección~\ref{sec:impl_esqueleto}.

\subsection{Clip de referencia}

El archivo de cada glosa funciona como el negativo de la seña. Para cada fotograma del video de origen conserva tres cosas: la posición de los puntos en píxeles de esa imagen, la cantidad de fotogramas por segundo y una confianza entre 0 y 1. La confianza dice cuán seguro está el estimador de haber visto ese punto. Esos datos quedan a la resolución del video original, sin achicarlos de antemano.

La escala al tamaño del lienzo se aplica en el momento de dibujar, y no dentro del archivo. El clasificador, en cambio, guarda las coordenadas como fracción del ancho y del alto, porque a la red le alcanza la forma del gesto. Si ese mismo número fraccionario se archivara para el avatar, el visor lo leería como una cantidad de píxeles: un valor $0{,}5$ pasaría a significar medio píxel, y el cuerpo quedaría comprimido contra la esquina del cuadro. El diccionario animable archiva por eso píxeles de imagen. La reducción a un tamaño común ocurre después, cuando ya se conoce el recuadro en el que la seña se va a mostrar.


El archivo crudo tampoco se recorta. Un video de diccionario muestra a la persona desde antes de la seña: levanta la mano, realiza el gesto y vuelve a una postura de reposo. El significado está en el tramo del medio, que aquí se llama núcleo. El levantamiento inicial es la preparación y el regreso al reposo es la retracción. El instante en que empieza el núcleo y el instante en que termina se anotan en un índice aparte, con un tiempo de entrada y un tiempo de salida por glosa. La reproducción lee ese índice. El negativo queda entero, de modo que un cambio del recorte no obliga a estimar el video otra vez.

Al dibujar no entran todos los puntos que el estimador produjo. MediaPipe Holistic marca 33 puntos de cuerpo, 468 de cara y 21 en cada mano. El diseño conserva el torso, de la cintura hacia arriba, y las dos manos completas. Las piernas se apagan, porque en este recuadro la seña se lee de medio cuerpo. Dentro de los 33 puntos de cuerpo también hay una cara gruesa y unas puntas de dedos muy simples. Esos puntos se apagan en el dibujo. Si se trazan junto con los 21 puntos de cada mano y con el contorno fino de la cara, el espectador ve una mano encima de la otra y una cara encima de la otra.

La posición del cuerpo expresada en metros, que el estimador entrega además de la posición en la imagen, se deja de lado por una razón distinta. Esa posición en metros no incluye las 21 articulaciones de la mano. Usarla como si fuera el esqueleto completo deja las manos pegadas al pecho aunque en el video estén junto a la cara. La lectura del avatar es frontal, y lo que importa es el lugar de cada punto en el plano de la imagen. La profundidad de cada mano, además, está medida respecto de su propia muñeca y en una escala mucho menor que el desplazamiento horizontal y vertical. Es un dato interno del estimador, y no una distancia en metros hasta la cámara.


\subsection{Preparación y empalme}

Dos archivos recién descargados todavía no pueden reproducirse uno detrás del otro. Cada video trae su propia velocidad, su propio encuadre y sus propios huecos allí donde el estimador perdió una mano. El diseño obliga a que todo clip atraviese la misma preparación, siempre en el mismo orden. El orden importa porque cada paso deja el material en condiciones para el siguiente: un hueco que no se rellena antes del empalme se convierte, en el empalme, en un salto.

\paragraph{Cadencia común.}
Los videos de origen no comparten necesariamente la misma cantidad de fotogramas por segundo. Si un clip trae 24 y el siguiente trae 30, una espera de seis fotogramas dura más en el primero que en el segundo, y el empalme pierde la noción de tiempo. Antes de cualquier otro retoque, todos los clips se llevan a una misma tasa. A partir de ahí, un fotograma significa la misma duración en todas las señas.

\paragraph{Recorte del núcleo.}
La reproducción usa el índice de entrada y de salida, y deja fuera la preparación y la retracción. Si se mostrara el video completo, quien mira vería a la persona llegar a la seña y retirarse de ella, y el final de una glosa se confundiría con el comienzo de la otra. El núcleo, que es el tramo donde la configuración y el movimiento distinguen esa seña de las demás, es lo único que pasa a la secuencia visible.

\paragraph{Relleno de huecos breves.}
Cuando el estimador deja de ver una mano durante unos pocos fotogramas, suele escribir el punto en el origen, es decir en la coordenada $(0, 0)$, o dejar un vacío entre dos fotogramas válidos. Si ese vacío es corto y está rodeado de fotogramas en los que la mano sí se vio, el diseño completa la posición con los vecinos válidos. El hueco largo, en el que la mano realmente no participa de la seña, no se inventa: lo resuelve el paso siguiente.

\paragraph{Mano en espera.}
Hay señas de una sola mano. En esos clips, el estimador puede copiar la mano visible sobre el lugar de la mano ausente, de modo que las dos manos quedan casi en el mismo sitio, o puede dejar la mano ausente en el origen. El diseño distingue esos dos casos de una mano que sí participa. Si las dos manos coinciden en una seña de una sola mano, la copia se retira. La mano que no participa se estaciona al costado del torso, y el antebrazo se dibuja como una línea continua desde el codo hasta la muñeca, para que la muñeca no aparezca quebrada en un ángulo que la seña no tiene. Esa postura de costado es la espera de la mano, y es distinta de la espera de todo el cuerpo, que se describe más abajo.

\paragraph{Suavizado.}
Aunque la mano esté bien seguida, la posición estimada tiembla de un fotograma al siguiente. El diseño mezcla cada medición nueva con la pose ya aceptada en el fotograma anterior, por un promedio exponencial. Un peso alto sigue al estimador con mayor fidelidad y deja pasar el temblor. Un peso bajo alisa de más y retrasa el gesto respecto del núcleo. El peso es el mismo para todo el diccionario. Su valor, y la ecuación, están en la Sección~\ref{sec:impl_esqueleto}.

\paragraph{Lienzo común.}
Los videos no comparten resolución ni tamaño de cuerpo. Si cada clip se dibujara a su escala original, al pasar de una glosa a la otra el esqueleto cambiaría de tamaño en medio de la frase. Después del suavizado, el cuerpo se lleva a un recuadro común y se escala hasta un mismo ancho de hombros. Dos señantes de distinta talla, y dos tomas de distinto encuadre, ocupan entonces el mismo espacio, y el empalme no produce un salto de escala.

Recién después de esos seis pasos el clip está en condiciones de unirse a otro. El dato que más daño hace, si se le deja llegar, es la coordenada nula. Supóngase que el último fotograma útil de \texttt{HOLA} tiene la muñeca en un punto interior del cuadro, y que el primero de \texttt{MAMÁ} trae esa misma muñeca en $(0, 0)$ porque el estimador no la vio. Una interpolación de unos pocos fotogramas entre esos dos puntos dibuja la mano en viaje hacia la esquina, como si la seña la llevara hasta allí. El relleno y la mano en espera existen para que el empalme reciba solo coordenadas válidas.
Entre el último fotograma del núcleo de una glosa y el primero del núcleo de la siguiente, el diseño inserta un puente: una serie breve de fotogramas nuevos, que no estaban en ninguno de los dos archivos. La serie es más larga cuando las muñecas de un extremo y del otro quedan más lejos, para que el recorrido no sea un salto. La curva con la que se reparte ese recorrido se llama \textit{smoothstep}: el punto sale despacio, acelera en la mitad y llega despacio, de modo que la unión no arranca con un tirón. Cada punto viaja en línea recta, en el plano de la imagen, desde donde está hasta donde tiene que estar. El archivo no contiene ángulos de huesos, y por eso el puente no gira una articulación: traslada los puntos.

Ese puente es un recurso gráfico. En la lengua, la coarticulación es otra cosa: el final de una seña ya anticipa la configuración de la siguiente. El puente no modela ese anticipo. Toma la pose final de un núcleo ya recortado y la pose inicial del núcleo siguiente, y cubre el espacio que hay entre ambas para que el ojo no vea un corte. El interior de cada seña llega al empalme tal como quedó después de la preparación.
La preparación de un clip tiene un costo de cálculo, y la descarga del archivo tiene un costo de red. Una vez que un clip quedó preparado, el diseño lo conserva en una caché en memoria. Si la misma glosa vuelve a aparecer en el enunciado, o en un enunciado posterior de la misma sesión, se reutiliza esa copia: no se vuelve a preparar y no se vuelve a pedir por la red. Además, la secuencia visible tiene un techo de 24 glosas y de 1200 fotogramas. Por encima de ese techo, un texto largo ocuparía el turno de la llamada con una animación que no termina. El diseño corta la secuencia en el techo. El Cuadro~\ref{tab:salvaguardas_pose} resume estas previsiones. El detalle de cada una está en los párrafos anteriores.
\begin{table}[H]
    \centering
    \small
    \caption{Previsiones del diseño frente a fallos del estimador y del léxico.}
    \label{tab:salvaguardas_pose}
    \begin{tabularx}{\textwidth}{@{}>{\raggedright\arraybackslash}p{0.42\textwidth}
                                 >{\raggedright\arraybackslash}X@{}}
        \toprule
        \textbf{Situación} & \textbf{Previsión} \\
        \midrule
        La mano desaparece durante unos pocos fotogramas y el estimador escribe $(0, 0)$. &
        Completa el hueco con los fotogramas vecinos en los que la mano sí se vio. \\
        \addlinespace[2pt]
        En una seña de una sola mano, las dos manos quedan en el mismo lugar. &
        Retira la copia y estaciona ese brazo al costado del torso. \\
        \addlinespace[2pt]
        La confianza de un punto es demasiado baja. &
        Omite ese punto en el trazo. \\
        \addlinespace[2pt]
        La posición estimada tiembla de un fotograma al otro. &
        Mezcla la medición nueva con la pose anterior, con un peso fijo. \\
        \addlinespace[2pt]
        La muñeca final de una glosa queda lejos de la muñeca inicial de la siguiente. &
        Inserta un puente breve de fotogramas nuevos entre los dos núcleos, sin alterar el interior de cada seña. \\
        \addlinespace[2pt]
        La glosa no tiene archivo. &
        La deletrea con los clips del alfabeto manual. \\
        \addlinespace[2pt]
        Tampoco están los archivos de las letras. &
        Deja el esqueleto en la postura de espera y avisa en la interfaz. \\
        \addlinespace[2pt]
        El enunciado supera 24 glosas o 1200 fotogramas. &
        Corta la secuencia en ese techo. \\
        \bottomrule
    \end{tabularx}
\end{table}

\subsection{Presentación en la videollamada}
\label{sec:presentacion_llamada}

Quien está del otro lado de la llamada ve la imagen que el navegador entrega como cámara. El diseño sustituye esa imagen por un lienzo partido en dos mitades de la misma altura. La mitad izquierda muestra la cámara de la persona, recortada para llenar su mitad sin estirar la figura. La mitad derecha muestra el esqueleto. Las dos mitades comparten la altura para que el intérprete y la persona queden a la misma escala.

Mientras no llega una lista de glosas, el esqueleto permanece en espera y no adelanta ninguna seña. La postura de espera se toma de un tramo en el que la persona todavía no empezó el gesto, y se escala al mismo ancho de hombros del lienzo. El antebrazo queda casi horizontal, levemente hacia abajo, y la mano sigue esa línea. La cara de la espera es un óvalo con cejas, ojos, nariz y boca, suficiente para situar la cabeza respecto de las manos.

Cuando la lista de glosas ya está resuelta y los fotogramas están empalmados, la reproducción no avanza un fotograma del clip por cada imagen que publica la llamada. La llamada puede publicar 24 imágenes por segundo mientras el clip quedó preparado a otra tasa. Si se avanzara al ritmo de la llamada, la seña saldría más lenta o más rápida que el video de origen. El diseño mide el tiempo transcurrido y elige el fotograma del clip que corresponde a ese tiempo.

Esa reproducción tiene tres tramos. En la entrada, un tramo breve a la tasa del clip, la postura de espera se transforma en el primer fotograma de la seña, y el contorno de cara de la seña aparece de a poco. En el tramo central corre la secuencia ya preparada, con los puentes entre glosas ya incluidos en la lista de fotogramas. En la salida, de la misma duración que la entrada, el último fotograma vuelve a la espera, el contorno de la seña se apaga y queda la cara simple de la espera. La curva de la entrada y de la salida es la misma \textit{smoothstep} del puente, de modo que el esqueleto no aparece ni desaparece de golpe. Al terminar, el clip se suelta y el esqueleto queda otra vez en espera, listo para la frase siguiente.

Hacia el lienzo no viajan los cientos de puntos del archivo. Viaja solo lo que hace falta para leer la seña en una mitad de cuadro: ocho puntos de pose, que son los dos hombros, los dos codos, las dos muñecas y las dos caderas; los 21 puntos de cada mano; y un contorno de cara formado por el óvalo, las cejas, los ojos, la nariz y los labios. En el recuadro pequeño de la llamada, esa malla se ve como una mancha densa y aporta poco a la configuración de los dedos, que es lo que la persona sorda necesita distinguir. Un punto de confianza baja se envía como ausente. El trazo lo saltea, en lugar de unir el punto anterior con una posición en la que el estimador casi no cree.
Cada dedo lleva un color fijo, y ese color es el mismo en las dos manos: pulgar rojo, índice naranja, medio amarillo, anular verde y meñique azul. La palma es blanca. El color permite seguir un dedo cuando las manos se cruzan delante del cuerpo y, sin él, los trazos de una y otra mano se confundirían.
La orientación también está fijada. En una toma de frente, la mano derecha de la persona aparece del lado izquierdo de la imagen, porque la cámara la ve de frente. El lienzo respeta esa orientación y no da vuelta el esqueleto para que la mano hábil quede del lado que el espectador llamaría derecho. La miniatura propia de la videollamada, en cambio, suele mostrarse espejada, como un reflejo. El diseño exige que esa persona y quien está del otro lado vean lo mismo: la cámara en la mitad izquierda y el esqueleto, sin invertir, en la mitad derecha. El mecanismo que deshace el espejo solo en la vista local queda en la Sección~\ref{sec:impl_esqueleto}.

\subsection{Contingencia léxica}
\label{sec:contingencia_lexica}

El nombre con el que se pide cada archivo sigue las reglas de la Sección~\ref{sec:api-avatar}: mayúsculas, sin tildes y sin espacios. Aquí el problema es otro. Una lista de glosas puede nombrar una seña con una palabra o con varias, y también puede nombrar algo que el diccionario no tiene.

La búsqueda prueba primero el grupo más largo. Toma tres glosas seguidas y pregunta si existe un archivo para ese grupo. Si no existe, prueba las dos primeras. Si tampoco existe, prueba una sola. El motivo es que una expresión como \texttt{BUENAS TARDES} puede ser una única seña, con un único archivo, y partirla de entrada en \texttt{BUENAS} y \texttt{TARDES} mostraría dos gestos distintos del que corresponde. Un dígito aislado, como el carácter \texttt{1}, se busca por el nombre de ese número, \texttt{UNO}, porque el diccionario nombra el número y no el signo con el que se escribe.

Cuando ni el grupo ni la palabra tienen archivo, el diseño deletrea la palabra con los clips del alfabeto manual. Así, un nombre propio como \texttt{FRANCISCO}, si no tiene seña propia, se muestra como la secuencia \texttt{F}, \texttt{R}, \texttt{A}, \texttt{N}, \texttt{C}, \texttt{I}, \texttt{S}, \texttt{C}, \texttt{O}. En ese alfabeto, \texttt{CH} es una sola letra, y no la letra \texttt{C} seguida de la letra \texttt{H}. \texttt{Ñ} también es una sola letra. Esas letras viajan juntas en \texttt{GET /fingerspell}, como fija la Sección~\ref{sec:ausencia_archivo}. Esta regla pertenece al deletreo. Es independiente de la regla del nombre de archivo, que solo decide cómo se escribe la cadena en el pedido.

Si falta el diccionario entero, faltan también las letras. En ese caso no hay con qué deletrear. El esqueleto permanece en la postura de espera, y la interfaz informa que el diccionario no está disponible. La palabra no se omite sin aviso, y tampoco se reemplaza por un movimiento aproximado: lo único que el lienzo puede mostrar es un clip que existe.

\section{Diseño de la extensión de videollamada}
\label{sec:diseno_extension}

La extensión une la videollamada con el motor local. El requerimiento RF-08 pide esa unión dentro de la pestaña de la llamada. La Figura~\ref{fig:arquitectura_video_ilsa} fija el reparto de cada bloque. Su construcción queda en el Capítulo~6.

\begin{figure}[!ht]
\centering
\resizebox{1\textwidth}{!}{%
\begin{circuitikz}
\tikzstyle{every node}=[font=\fontsize{8.0pt}{10.4pt}\selectfont]

% Marco principal
\draw [line width=0.7pt] (-3.375,15.125) rectangle (9.125,6.875);

% Cámara web
\draw (-1.875,16.75) rectangle (1.375,16);
\node[anchor=center, align=center, inner xsep=0.080cm, inner ysep=0.085cm, rounded corners=0.020cm, font=\fontsize{10.8pt}{14.1pt}\selectfont] at (-0.25,16.375) {Cámara web};

% Google Meet principal
\draw [line width=0.7pt] (-3.125,14.875) rectangle (0.875,7.125);
\node[anchor=north west, align=left, inner xsep=0.080cm, inner ysep=0.085cm, rounded corners=0.020cm, font=\fontsize{10.8pt}{14.1pt}\selectfont] at (-2.975,14.725) {Google Meet};

% HUD
\draw (2.3,7.875) rectangle (5.05,7.25);
\node[anchor=center, align=center, inner xsep=0.080cm, inner ysep=0.085cm, rounded corners=0.020cm, font=\fontsize{10.8pt}{14.1pt}\selectfont] at (3.675,7.5625) {indicador};

% Offscreen
\draw [line width=0.7pt] (2.5,14.75) rectangle (7,10.625);
\node[anchor=north east, align=right, inner xsep=0.080cm, inner ysep=0.085cm, rounded corners=0.020cm, font=\fontsize{10.8pt}{14.1pt}\selectfont] at (6.85,14.6) {coordinación};

% ILSA (motor local)
\draw (-3.25,4.75) rectangle (9.125,4);
\node[anchor=center, align=center, inner xsep=0.080cm, inner ysep=0.085cm, rounded corners=0.020cm, font=\fontsize{10.8pt}{14.1pt}\selectfont] at (2.9375,4.375) {ILSA (motor local)};

% Cajas dentro de Google Meet
\draw (-2.875,12.375) rectangle (0.125,11.25);
\node[anchor=center, align=center, inner xsep=0.080cm, inner ysep=0.085cm, rounded corners=0.020cm, font=\fontsize{10.8pt}{14.1pt}\selectfont] at (-1.375,11.8125) {video oculto \\ (flujo real)};

\draw (-2.875,14) rectangle (0.125,12.875);
\node[anchor=center, align=center, inner xsep=0.080cm, inner ysep=0.085cm, rounded corners=0.020cm, font=\fontsize{10.8pt}{14.1pt}\selectfont] at (-1.375,13.4375) {toma de la\\ cámara};

\draw (-2.875,10.25) rectangle (0.625,9.125);
\node[anchor=center, align=center, inner xsep=0.080cm, inner ysep=0.085cm, rounded corners=0.020cm, font=\fontsize{10.8pt}{14.1pt}\selectfont] at (-1.125,9.6875) {lienzo \\ (imagen y subtítulo)};

\draw (-2.875,8.5) rectangle (0.125,7.375);
\node[anchor=center, align=center, inner xsep=0.080cm, inner ysep=0.085cm, rounded corners=0.020cm, font=\fontsize{10.8pt}{14.1pt}\selectfont] at (-1.375,7.9375) {Google Meet \\ (videollamada)};

% Cajas dentro de Offscreen
\draw (2.75,14) rectangle (6.75,12.75);
\node[anchor=center, align=center, inner xsep=0.080cm, inner ysep=0.085cm, rounded corners=0.020cm, font=\fontsize{10.8pt}{14.1pt}\selectfont] at (4.75,13.375) {controlador \\ (fotogramas reducidos)};

\draw (2.75,11.75) rectangle (6.75,10.875);
\node[anchor=center, align=center, inner xsep=0.080cm, inner ysep=0.085cm, rounded corners=0.020cm, font=\fontsize{10.8pt}{14.1pt}\selectfont] at (4.75,11.3125) {temporizador};

% Micrófono o chat
\draw (4.625,6.625) rectangle (9.125,5.875);
\node[anchor=center, align=center, inner xsep=0.080cm, inner ysep=0.085cm, rounded corners=0.020cm, font=\fontsize{10.8pt}{14.1pt}\selectfont] at (6.875,6.25) {Micrófono o chat};

% Sandbox
\draw [line width=0.7pt] (4.25,10.25) rectangle (8.75,8.375);
\node[anchor=north west, align=left, inner xsep=0.080cm, inner ysep=0.085cm, rounded corners=0.020cm, font=\fontsize{10.8pt}{14.1pt}\selectfont] at (4.4,10.1) {visión};

\draw (4.5,9.5) rectangle (8.5,8.625);
\node[anchor=center, align=center, inner xsep=0.080cm, inner ysep=0.085cm, rounded corners=0.020cm, font=\fontsize{10.8pt}{14.1pt}\selectfont] at (6.5,9.0625) {extracción de puntos};

% Conexiones y flechas
\draw [-{Stealth[scale=1.5]}] (0,16) -- (0,14);
\draw [-{Stealth[scale=1.5]}] (-1.25,12.875) -- (-1.25,12.375);
\draw [-{Stealth[scale=1.5]}] (-1.375,11.25) -- (-1.375,10.25) node[pos=0.5, fill=white, fill opacity=1, text opacity=1, inner xsep=0.080cm, inner ysep=0.085cm, rounded corners=0.020cm]{cuadro a cuadro};
\draw [-{Stealth[scale=1.5]}] (-1.375,9.125) -- (-1.375,8.5);

\draw [-{Stealth[scale=1.5]}] (0.125,11.875) -- (1.625,11.875) -- (1.625,13.5) -- (2.75,13.5);
\node[fill=white, fill opacity=1, text opacity=1, inner xsep=0.080cm, inner ysep=0.085cm, rounded corners=0.020cm, font=\fontsize{8.0pt}{10.4pt}\selectfont] at (1.625,12.5) {JPEG};

\draw [-{Stealth[scale=1.5]}] (6.75,13.875) -- (8.25,13.875) -- (8.25,9.5);
\node[fill=white, fill opacity=1, text opacity=1, inner xsep=0.080cm, inner ysep=0.085cm, rounded corners=0.020cm, font=\fontsize{8.0pt}{10.4pt}\selectfont] at (8.25,12.5) {fotograma};

\draw [dash pattern=on 6pt off 2pt, -{Stealth[scale=1.5]}] (7.625,9.5) -- (7.625,13.25) -- (6.75,13.25);
\node[fill=white, fill opacity=1, text opacity=1, inner xsep=0.080cm, inner ysep=0.085cm, rounded corners=0.020cm, font=\fontsize{8.0pt}{10.4pt}\selectfont] at (7.625,11.875) {puntos};

\draw [-{Stealth[scale=1.5]}] (2.75,11.25) -- (1.625,11.25) -- (1.625,4.75);
\node[fill=white, fill opacity=1, text opacity=1, inner xsep=0.080cm, inner ysep=0.085cm, rounded corners=0.020cm, font=\fontsize{8.0pt}{10.4pt}\selectfont] at (1.625,6.5) {POST /sign};

\draw [-{Stealth[scale=1.5]}] (5,5.875) -- (5,4.75);
\node[fill=white, fill opacity=1, text opacity=1, inner xsep=0.080cm, inner ysep=0.085cm, rounded corners=0.020cm, font=\fontsize{8.0pt}{10.4pt}\selectfont] at (5,5.375) {POST /stt};

\draw [dash pattern=on 6pt off 2pt, -{Stealth[scale=1.5]}] (6.5,4.75) -- (6.5,5.875);
\node[fill=white, fill opacity=1, text opacity=1, inner xsep=0.080cm, inner ysep=0.085cm, rounded corners=0.020cm, font=\fontsize{8.0pt}{10.4pt}\selectfont] at (6.5,5.375) {oración};

\draw [-{Stealth[scale=1.5]}] (8.375,5.875) -- (8.375,4.75);
\node[fill=white, fill opacity=1, text opacity=1, inner xsep=0.080cm, inner ysep=0.085cm, rounded corners=0.020cm, font=\fontsize{8.0pt}{10.4pt}\selectfont] at (8.375,5.375) {POST /hearing};

\draw [-{Stealth[scale=1.5]}] (3.375,4.75) -- (3.375,7.25);
\node[fill=white, fill opacity=1, text opacity=1, inner xsep=0.080cm, inner ysep=0.085cm, rounded corners=0.020cm, font=\fontsize{8.0pt}{10.4pt}\selectfont] at (3.375,6) {estado / glosas};

\draw [dash pattern=on 6pt off 2pt, -{Stealth[scale=1.5]}] (-2.875,4.75) -- (-2.875,5.75) -- (-4.75,5.75) -- (-4.75,9.75) -- (-2.875,9.75);
\node[fill=white, fill opacity=1, text opacity=1, inner xsep=0.080cm, inner ysep=0.085cm, rounded corners=0.020cm, font=\fontsize{8.0pt}{10.4pt}\selectfont] at (-3.875,5.75) {animación};

\draw [dash pattern=on 6pt off 2pt, -{Stealth[scale=1.5]}] (-2.375,4.75) -- (-2.375,6.375) -- (-4.125,6.375) -- (-4.125,9.375) -- (-2.875,9.375);
\node[fill=white, fill opacity=1, text opacity=1, inner xsep=0.080cm, inner ysep=0.085cm, rounded corners=0.020cm, font=\fontsize{8.0pt}{10.4pt}\selectfont] at (-3.25,6.375) {subtítulos};

\draw [-{Stealth[scale=1.5]}] (3.625,12.75) -- (3.625,11.75);
\node[fill=white, fill opacity=1, text opacity=1, inner xsep=0.080cm, inner ysep=0.085cm, rounded corners=0.020cm, font=\fontsize{8.0pt}{10.4pt}\selectfont] at (3.625,12.375) {puntos};

\end{circuitikz}%
}
\caption{Arquitectura de la extensión de videollamada. El motor local queda como un solo componente.}
\label{fig:arquitectura_video_ilsa}
\end{figure}

\subsection{Frontera}
La extensión es la única pieza que conoce Google Meet. No clasifica una seña, no arma la oración y no guarda el enunciado. Esas tareas pertenecen a ILSA. El búfer de glosas vive en el motor, como ya fija el Cuadro~\ref{tab:api-ilsa-local}. Toda llamada de la extensión usa la dirección local de ese cuadro.

En modo señante el temporizador envía a \texttt{POST /sign} los puntos de la seña ya cerrada. ILSA devuelve la oración. Esa oración llega al lienzo como subtítulo. El estado y las glosas llegan al indicador.

En modo oyente el micrófono entrega el audio a \texttt{POST /stt} y recibe la oración transcrita. Esa oración, o el texto ya escrito en el chat, entra por \texttt{POST /hearing}. ILSA devuelve la animación y los subtítulos al lienzo, y las glosas al indicador. La extensión no llama al traductor ni al servicio de archivos. Esa llamada sale de ILSA.

\subsection{Dos vistas de la misma llamada}
Quien está del otro lado solo ve la imagen que la llamada transmite. En modo señante, esa imagen lleva el subtítulo de la oración. En modo oyente, lleva el lienzo partido de la Sección~\ref{sec:presentacion_llamada}: la cámara en una mitad y el esqueleto en la otra, junto con los subtítulos.

Quien tiene la extensión ve, además, el indicador de la Figura~\ref{fig:arquitectura_video_ilsa}. El indicador informa si el motor responde, cuál fue la última glosa y si el turno sigue abierto. Ese aviso no entra en la imagen transmitida. Así no tapa al interlocutor y se cumple el requerimiento RNF-05.

\subsection{Aislamiento}
Dentro del navegador el diseño separa tres ámbitos. La Figura~\ref{fig:arquitectura_video_ilsa} los dibuja como tres bloques.

El ámbito de la página es el bloque Google Meet. La toma de la cámara recibe la cámara web. El video oculto conserva el flujo real. El lienzo recibe esa imagen cuadro a cuadro, y la llamada publica el lienzo. Este ámbito es el único que toca la videollamada.

El ámbito de coordinación agrupa al controlador y al temporizador. El controlador recibe del video oculto un JPEG ya reducido, lo entrega a la visión y recibe los puntos. Esos puntos pasan al temporizador. En modo señante, el temporizador habla con ILSA por \texttt{POST /sign}.

El ámbito de visión extrae los puntos. Recibe un fotograma y devuelve los 225 puntos de pose y manos. No conoce Meet ni la dirección del motor.

El micrófono o el chat quedan fuera de esos tres bloques. En modo oyente hablan con ILSA. Ni el controlador ni la visión entran en ese camino.

Entre los tres ámbitos no circula el video crudo. Hacia la red salen solo glosas u oraciones, y salen desde ILSA. Los puntos articulares permanecen en el equipo. Esa partición sostiene el requerimiento RNF-02.

\subsection{Fallo de un servicio}
La llamada continúa aunque un servicio no responda.

Si ILSA no responde, no hay traducción. La cámara de la llamada sigue como cámara de la persona. La extensión no publica un lienzo vacío en su lugar, y el indicador local muestra que el motor no está disponible.

Si el servicio de archivos no responde, vale la regla de la Sección~\ref{sec:contingencia_lexica}. El esqueleto queda en la postura de espera, la interfaz avisa, y las glosas siguen disponibles como texto.

\section{Patrones de diseño de software aplicados}

El diseño separa las responsabilidades del cliente local, de la extensión y del servicio remoto. Cada patrón acota un problema distinto: la política de admisión de glosas, la interfaz única hacia el motor, la conversión entre contratos y el modo de la sesión. El efecto de cada decisión queda en el componente dueño. La red no filtra rebotes, la extensión no llama a la nube y el servicio remoto no altera el estado de la sesión local. La referencia de estos patrones es \parencite{gamma1994design}.

\begin{itemize}
    \item \textbf{Estrategia (\textit{Strategy}), en el filtro de repetición.}
    El clasificador propone una glosa y una confianza. No decide si esa glosa entra al enunciado. Esa decisión queda en un objeto de política, independiente de la red. La política distingue tres categorías del léxico: los dígitos admiten repetición, las letras admiten hasta dos apariciones consecutivas y el resto del vocabulario rechaza la repetición inmediata. El acumulador de enunciado consulta la política y conserva el reloj de pausa Un cambio de la regla de repetición no altera el modelo ni el contrato remoto.

    \item \textbf{Fachada (\textit{Facade}), en la API local de ILSA.}
    La extensión no invoca por separado al clasificador, a la compuerta, al acumulador ni al cliente de la nube. El motor expone un único contrato HTTP en el bucle local: sesión, ingesta de la seña, aviso de actividad, cierre de enunciado, estado y diagnóstico. La fachada oculta el orden interno del pipeline y concentra el acceso del navegador.

    \item \textbf{Adaptador (\textit{Adapter}), en dos fronteras.}
    La extensión adapta los medios de Google Meet al contrato local. La cámara y el lienzo de la llamada se presentan como fotogramas y subtítulos, y el audio o el texto del oyente se presentan como una oración. En la otra frontera, ILSA adapta el enunciado cerrado al contrato del servicio remoto: una cadena de glosas en \texttt{/sordo} o en \texttt{/translate}, o una oración en \texttt{/hearing}. Con las glosas del canal inverso, ILSA arma el nombre de archivo y consulta \texttt{/tposes}. Si esa glosa no tiene archivo, las letras salen por \texttt{/fingerspell}. La respuesta remota vuelve a entrar por ILSA. Ni el traductor ni el servicio de archivos conocen Meet ni la extensión.

    \item \textbf{Estado (\textit{State}), en el modo de sesión.}
    La sesión admite dos modos excluyentes: señante y oyente. En modo señante el motor acepta señas y cierra enunciados hacia el canal directo. En modo oyente acepta el texto del interlocutor, consulta el canal inverso y pide los archivos de seña. La misma fachada cambia de comportamiento según el modo, sin una segunda extensión. El temporizador de pausa no es un estado del clasificador: pertenece al acumulador y solo declara el cierre del enunciado.
\end{itemize}

\section{Diseño conceptual de la arquitectura escalable de producción en la nube}

Las pruebas de este trabajo usan una instancia de Cloud Run por servicio. Esta sección no describe esa instancia. Describe la proyección para varias videollamadas a la vez, en coincidencia con la restricción RC-02. 

\begin{figure}[!ht]
\centering
\resizebox{1\textwidth}{!}{%
\begin{circuitikz}
\tikzstyle{every node}=[font=\fontsize{8.0pt}{10.4pt}\selectfont]

% Clientes concurrentes
\draw [line width=0.7pt] (3.125,12.875) rectangle (5.25,9.625);
\node[anchor=north, align=center, inner xsep=0.080cm, inner ysep=0.085cm, rounded corners=0.020cm, font=\fontsize{8.0pt}{10.4pt}\selectfont] at (4.1875,12.725) {Clientes \\ concurrentes};

\draw (3.375,11.875) rectangle (5,11.375);
\node[anchor=center, align=center, inner xsep=0.080cm, inner ysep=0.085cm, rounded corners=0.020cm, font=\fontsize{8.0pt}{10.4pt}\selectfont] at (4.1875,11.625) {Cliente A};

\draw (3.375,11.125) rectangle (5,10.625);
\node[anchor=center, align=center, inner xsep=0.080cm, inner ysep=0.085cm, rounded corners=0.020cm, font=\fontsize{8.0pt}{10.4pt}\selectfont] at (4.1875,10.875) {Cliente A};

\draw (3.375,10.375) rectangle (5,9.875);
\node[anchor=center, align=center, inner xsep=0.080cm, inner ysep=0.085cm, rounded corners=0.020cm, font=\fontsize{8.0pt}{10.4pt}\selectfont] at (4.1875,10.125) {Cliente A};

% Ingreso
\draw [line width=0.7pt] (6.5,11.875) rectangle (13.125,10.125);
\node[anchor=north west, align=left, inner xsep=0.080cm, inner ysep=0.085cm, rounded corners=0.020cm, font=\fontsize{8.0pt}{10.4pt}\selectfont] at (6.65,11.725) {Ingreso};

\draw (6.875,11.125) rectangle (9.625,10.625);
\node[anchor=center, align=center, inner xsep=0.080cm, inner ysep=0.085cm, rounded corners=0.020cm, font=\fontsize{8.0pt}{10.4pt}\selectfont] at (8.25,10.875) {DNS / CDN / WAF};

\draw (10.25,11.625) rectangle (12.75,10.375);
\node[anchor=center, align=center, inner xsep=0.080cm, inner ysep=0.085cm, rounded corners=0.020cm, font=\fontsize{8.0pt}{10.4pt}\selectfont] at (11.5,11) {\textit{Aplication load} \\ \textit{balancer / Ingress} \\ \textit{controller}};

% Servicio semántico
\draw [line width=0.7pt] (14,15.375) rectangle (23.875,11.5);
\node[anchor=north west, align=left, inner xsep=0.080cm, inner ysep=0.085cm, rounded corners=0.020cm, font=\fontsize{8.0pt}{10.4pt}\selectfont] at (14.15,15.225) {Servicio semántico (sin estado, cómputo)};

\draw (14.875,14.75) rectangle (17.375,11.75);
\node[anchor=north west, align=left, inner xsep=0.080cm, inner ysep=0.085cm, rounded corners=0.020cm, font=\fontsize{8.0pt}{10.4pt}\selectfont] at (15.025,14.6) {Réplicas idénticas};

\draw (15.25,14) rectangle (16.875,13.5);
\node[anchor=center, align=center, inner xsep=0.080cm, inner ysep=0.085cm, rounded corners=0.020cm, font=\fontsize{8.0pt}{10.4pt}\selectfont] at (16.0625,13.75) {Réplica 1};

\draw (15.25,13.25) rectangle (16.875,12.75);
\node[anchor=center, align=center, inner xsep=0.080cm, inner ysep=0.085cm, rounded corners=0.020cm, font=\fontsize{8.0pt}{10.4pt}\selectfont] at (16.0625,13) {Réplica 2};

\draw (15.25,12.5) rectangle (16.875,12);
\node[anchor=center, align=center, inner xsep=0.080cm, inner ysep=0.085cm, rounded corners=0.020cm, font=\fontsize{8.0pt}{10.4pt}\selectfont] at (16.0625,12.25) {Réplica N};

\draw (21.125,15) rectangle (23.375,13.625);
\node[anchor=north, align=center, inner xsep=0.080cm, inner ysep=0.085cm, rounded corners=0.020cm, font=\fontsize{8.0pt}{10.4pt}\selectfont] at (22.25,14.85) {Escalador \\ horizontal (cola \\ y latencia)};

\draw (21.125,13.125) rectangle (23.25,11.75);
\node[anchor=north, align=center, inner xsep=0.080cm, inner ysep=0.085cm, rounded corners=0.020cm, font=\fontsize{8.0pt}{10.4pt}\selectfont] at (22.1875,12.975) {Modelos en \\ volumen (solo \\ lectura)};

% Servicio de poses
\draw [line width=0.7pt] (14,10.125) rectangle (23.875,6.25);
\node[anchor=north west, align=left, inner xsep=0.080cm, inner ysep=0.085cm, rounded corners=0.020cm, font=\fontsize{8.0pt}{10.4pt}\selectfont] at (14.15,9.975) {Servicio de poses (sin estado, I/O)};

\draw (14.75,9.5) rectangle (17.25,6.5);
\node[anchor=north west, align=left, inner xsep=0.080cm, inner ysep=0.085cm, rounded corners=0.020cm, font=\fontsize{8.0pt}{10.4pt}\selectfont] at (14.9,9.35) {Réplicas idénticas};

\draw (15.125,8.75) rectangle (16.75,8.25);
\node[anchor=center, align=center, inner xsep=0.080cm, inner ysep=0.085cm, rounded corners=0.020cm, font=\fontsize{8.0pt}{10.4pt}\selectfont] at (15.9375,8.5) {Réplica 1};

\draw (15.125,8) rectangle (16.75,7.5);
\node[anchor=center, align=center, inner xsep=0.080cm, inner ysep=0.085cm, rounded corners=0.020cm, font=\fontsize{8.0pt}{10.4pt}\selectfont] at (15.9375,7.75) {Réplica 2};

\draw (15.125,7.25) rectangle (16.75,6.75);
\node[anchor=center, align=center, inner xsep=0.080cm, inner ysep=0.085cm, rounded corners=0.020cm, font=\fontsize{8.0pt}{10.4pt}\selectfont] at (15.9375,7) {Réplica N};

\draw (20.75,9.75) rectangle (23.375,8.375);
\node[anchor=north, align=center, inner xsep=0.080cm, inner ysep=0.085cm, rounded corners=0.020cm, font=\fontsize{8.0pt}{10.4pt}\selectfont] at (22.0625,9.6) {Escalador \\ horizontal (QPS y \\ ancho de banda)};

\draw (21.125,7.875) rectangle (23.25,6.5);
\node[anchor=north, align=center, inner xsep=0.080cm, inner ysep=0.085cm, rounded corners=0.020cm, font=\fontsize{8.0pt}{10.4pt}\selectfont] at (22.1875,7.725) {Carpeta de \\ poses (solo \\ lectura)};

% Observabilidad
\draw [line width=0.7pt] (25,11.625) rectangle (27.625,10.125);
\node[anchor=north, align=center, inner xsep=0.080cm, inner ysep=0.085cm, rounded corners=0.020cm, font=\fontsize{8.0pt}{10.4pt}\selectfont] at (26.3125,11.475) {Observabilidad};

\draw (25.5,11) rectangle (27.125,10.5);
\node[anchor=center, align=center, inner xsep=0.080cm, inner ysep=0.085cm, rounded corners=0.020cm, font=\fontsize{8.0pt}{10.4pt}\selectfont] at (26.3125,10.75) {Métricas};

% Notas informativas
\draw [draw opacity=0, fill={rgb,255:red,231; green,231; blue,231}, fill opacity=1] (10.25,12.125) rectangle (12.75,13.125);
\node[anchor=north, align=center, inner xsep=0.080cm, inner ysep=0.085cm, rounded corners=0.020cm, font=\fontsize{8.0pt}{10.4pt}\selectfont] at (11.5,12.975) {POST /sign \\ POST /hearing};

\draw [draw opacity=0, fill={rgb,255:red,231; green,231; blue,231}, fill opacity=1] (9.25,8.25) rectangle (13.125,9.25);
\node[anchor=north, align=center, inner xsep=0.080cm, inner ysep=0.085cm, rounded corners=0.020cm, font=\fontsize{8.0pt}{10.4pt}\selectfont] at (11.1875,9.1) {GET /tsposes/\{archivo\} \\ GET /fingerspell?letters=...};

\draw [draw opacity=0, fill={rgb,255:red,231; green,231; blue,231}, fill opacity=1] (18.125,11.625) rectangle (20.75,13.375);
\node[anchor=north, align=center, inner xsep=0.080cm, inner ysep=0.085cm, rounded corners=0.020cm, font=\fontsize{8.0pt}{10.4pt}\selectfont] at (19.4375,13.225) {Montaje \\ compartido, el peso \\ vive en RAM de \\ cada réplica};

\draw [draw opacity=0, fill={rgb,255:red,231; green,231; blue,231}, fill opacity=1] (17.875,6.5) rectangle (20.5,7.5);
\node[anchor=north, align=center, inner xsep=0.080cm, inner ysep=0.085cm, rounded corners=0.020cm, font=\fontsize{8.0pt}{10.4pt}\selectfont] at (19.1875,7.35) {Misma carpeta en \\ cada réplica};

% Conexiones y flechas
\draw [-{Stealth[scale=1.5]}] (5,10.875) -- (6.875,10.875) node[pos=0.33, fill=white, fill opacity=1, text opacity=1, inner xsep=0.080cm, inner ysep=0.085cm, rounded corners=0.020cm]{HTTPS};

\draw [-{Stealth[scale=1.5]}] (5,11.625) -- (6.125,11.625) -- (6.125,11.125) -- (6.875,11.125);
\node[fill=white, fill opacity=1, text opacity=1, inner xsep=0.080cm, inner ysep=0.085cm, rounded corners=0.020cm, font=\fontsize{8.0pt}{10.4pt}\selectfont] at (5.625,11.625) {HTTPS};

\draw [-{Stealth[scale=1.5]}] (5,10.125) -- (6.125,10.125) -- (6.125,10.625) -- (6.875,10.625);
\node[fill=white, fill opacity=1, text opacity=1, inner xsep=0.080cm, inner ysep=0.085cm, rounded corners=0.020cm, font=\fontsize{8.0pt}{10.4pt}\selectfont] at (5.75,10.125) {HTTPS};

\draw [-{Stealth[scale=1.5]}] (9.625,10.875) -- (10.25,10.875);

\draw [-{Stealth[scale=1.5]}] (11.625,11.625) -- (11.625,12.25) -- (15.25,12.25);
\draw [-{Stealth[scale=1.5]}] (11.625,11.625) -- (11.625,13) -- (15.25,13);
\draw [-{Stealth[scale=1.5]}] (11.625,11.625) -- (11.625,13.625) -- (15.25,13.625);

\draw [-{Stealth[scale=1.5]}] (11.625,10.375) -- (11.625,8.625) -- (15.125,8.625);
\draw [-{Stealth[scale=1.5]}] (11.625,10.375) -- (11.625,7.875) -- (15.125,7.875);
\draw [-{Stealth[scale=1.5]}] (11.625,10.375) -- (11.625,7) -- (15.125,7);

\draw [-{Stealth[scale=1.5]}] (20.75,8.625) -- (17.25,8.625) node[pos=0.5, fill=white, fill opacity=1, text opacity=1, inner xsep=0.080cm, inner ysep=0.085cm, rounded corners=0.020cm]{suma o baja réplicas};
\draw [-{Stealth[scale=1.5]}] (21.125,7) -- (17.25,7);
\draw [-{Stealth[scale=1.5]}] (21.125,12.375) -- (17.375,12.375);
\draw [-{Stealth[scale=1.5]}] (21.125,14.25) -- (17.375,14.25) node[pos=0.5, fill=white, fill opacity=1, text opacity=1, inner xsep=0.080cm, inner ysep=0.085cm, rounded corners=0.020cm]{suma o baja réplicas};

\draw [-{Stealth[scale=1.5]}] (17.25,9.125) -- (18.75,9.125) -- (18.75,10.625) -- (25.5,10.625);
\node[fill=white, fill opacity=1, text opacity=1, inner xsep=0.080cm, inner ysep=0.085cm, rounded corners=0.020cm, font=\fontsize{8.0pt}{10.4pt}\selectfont] at (22.125,10.625) {latencia, ocupación, CPU};

\draw [-{Stealth[scale=1.5]}] (16.875,11.75) -- (16.875,11) -- (25.5,11);
\node[fill=white, fill opacity=1, text opacity=1, inner xsep=0.080cm, inner ysep=0.085cm, rounded corners=0.020cm, font=\fontsize{8.0pt}{10.4pt}\selectfont] at (18.625,11) {lecturas, errores, bytes};

\draw [-{Stealth[scale=1.5]}] (26.375,11.625) -- (26.375,14.25) -- (23.375,14.25);
\draw [-{Stealth[scale=1.5]}] (26.375,10.125) -- (26.375,9.125) -- (23.375,9.125);

\end{circuitikz}%
}
\caption{Diagrama de la arquitectura del servidor, balanceo de carga y escalabilidad de servicios sin estado.}
\label{fig:arquitectura_servidor_escalabilidad}
\end{figure}

La proyección tiene dos servicios detrás de un mismo ingreso. El traductor genera texto. El servicio de archivos entrega un clip de seña. No comparten réplica ni cola.

\subsection{Ingreso}
Las computadoras cliente salen por HTTPS. El ingreso resuelve el nombre, presenta el certificado TLS y filtra tráfico elemental. Detrás opera un balanceador.
\begin{itemize}
    \item \texttt{POST /sign} y \texttt{POST /hearing} van al traductor. El criterio es el menor número de conexiones activas, para no cargar una réplica que ya genera un enunciado largo.
    \item \texttt{GET /tposes/\{glosa\}.pose}, \texttt{GET /fingerspell} y la salud de ese servicio van a las réplicas de archivos.
\end{itemize}

\subsection{Traductor}
El traductor viaja en una imagen de contenedor. Cada réplica monta los binarios GGUF Q4\_K\_M desde un volumen de solo lectura, así un arranque nuevo no vuelve a bajar los modelos. Un escalador horizontal mira la cola y el tiempo de generación. Si los turnos simultáneos superan a las réplicas activas, el orquestador abre contenedores nuevos. Por la cuantización a 4 bits, cada réplica cabe en menos de 2\,GB de RAM y corre en procesador, sin placa de video.

\subsection{Archivos de seña}
Una petición de pose es una lectura, no una generación. El contrato es el de la Sección~\ref{sec:api-avatar}. Las réplicas no guardan estado y leen la misma carpeta. El escalador de este grupo mira la cantidad de lecturas, no la cola del traductor. Un paso posterior puede pasar la carpeta a un almacén de objetos y dejar una red de contenido delante de las glosas repetidas. La ruta de salud queda aparte de la ruta del archivo. La caché del cliente sigue como primer recurso.

\subsection{Monitoreo}
Un recolector de métricas observa la latencia por oración, la memoria de cada contenedor del traductor, los errores de red y los errores de lectura del servicio de archivos. Esas series ajustan los dos escaladores.

\subsection{Qué queda fuera de las pruebas}
El prototipo no crea el balanceador, ni los escaladores, ni el volumen compartido. ILSA consulta una instancia del traductor, con los dos canales en un solo proceso, y una instancia del servicio de archivos. Esas dos publicaciones se describen en el Capítulo~6. La carga concurrente y el auto-escalado no entran en la validación.
